% Protocol version: everything up to and including the Method section was
% written and committed before any literature search was run.  Later
% changes to those sections are reported in the section on deviations.

\section{Introduction}\label{sec:intro}

Multi-factor authentication is commonly recommended as the way to make
user authentication secure, and in places required: the United States
government requires it of its agencies' staff\parencite{OMB2022ZeroTrust},
and its cybersecurity agency writes that
\textcquote[p.~1]{CISA2022PhishingResistant}{any form of MFA is better
than no MFA}.
This report asks whether the literature supports two stronger readings of
that recommendation: that multi-factor authentication is \emph{sufficient}
for secure authentication of users, and that it is \emph{necessary}.

Neither reading can be tested as it stands.
Whether something is sufficient or necessary for security depends on what
counts as multi-factor authentication, on what the attacker can do, and on
what we compare with.
In particular, passwordless authentication replaces the password with a
cryptographic key.
If that works, the benefit usually credited to having \emph{several}
factors may instead come from having \emph{one strong} factor, and the
password may contribute nothing once the key is there.
We therefore fix our terms (\cref{sec:terms}), the threats
(\cref{sec:threats}) and the meaning of necessary and sufficient
(\cref{sec:logic}) before posing the questions (\cref{sec:questions}).

The report is also meant as an example of how to explore the literature
on a question of this kind.
For that reason the questions, the predicted answers and the method were
written down before the search, and the report says where the search
departed from them (\cref{sec:deviations}).

\subsection{Terms}\label{sec:terms}

The following are working definitions, fixed before the search.
How the standards define the same terms is one of the things the search
is to establish; \cref{sec:results} reports any difference.
\begin{description}
  \item[Authentication factor] A kind of evidence of identity: something
    the user knows, something the user has, or something the user is.
  \item[Authenticator] The concrete thing that provides the evidence: a
    password, a one-time code sent by text message, a code generated by an
    app, an approval of a push notification, a cryptographic key held in a
    hardware token or a phone.
  \item[Multi-factor authentication] Authentication that requires evidence
    of at least two different kinds of factor.
    Two-factor authentication is the special case of exactly two.
    Two steps of the same kind, such as a password and a security
    question, are not multi-factor authentication in this sense.
  \item[Passwordless authentication] Authentication that uses no secret
    memorised by the user and shared with the service.
    It says what is absent, not how many factors remain: it can be
    single-factor or multi-factor.
  \item[User verification] A local check, by a PIN or a biometric, that
    unlocks a cryptographic key on the user's device.
    The PIN or biometric never reaches the service.
\end{description}

To separate the number of factors from their strength we compare the
cases in \cref{tab:cases}.
\begin{table}
  \begin{sidecaption}{%
    The cases compared.
    Case \case{B} is passwordless but multi-factor; case \case{C} is
    passwordless and single-factor.
    Comparing \case{A3} with \case{B} and \case{C} shows what the password
    contributes beside a key; comparing \case{P+} with \case{A1}--\case{A3}
    shows what a second factor contributes beside a strong password.
  }[tab:cases]
    \begin{tabular}{@{}lll@{}}
      \toprule
      Case & What the user presents & Factors \\
      \midrule
      \case{P}  & Password, chosen by the user & one \\
      \case{P+} & Password, random and unique (password manager) & one \\
      \case{A1} & Password and a code (text message, app) & two \\
      \case{A2} & Password and a push approval & two \\
      \case{A3} & Password and a cryptographic key & two \\
      \case{B}  & Cryptographic key with user verification & two \\
      \case{C}  & Cryptographic key alone & one \\
      \bottomrule
    \end{tabular}
  \end{sidecaption}
\end{table}

\subsection{Threat model}\label{sec:threats}

We take \enquote{secure} to mean that an attacker with a stated
capability cannot authenticate as the user.
\Cref{tab:threats} lists the capabilities considered.
The remote threats can be carried out at scale against many users; the
local threats need access to one user's device.
Threats \threat{R7} and \threat{R8} act around the authenticator and not
on it.
We include them because the question concerns the authentication of
users, and a user's account is compromised whichever way the attacker
gets in.
\begin{table}
  \begin{sidecaption}{%
    The threat model: what the attacker is assumed able to do.
    R is remote, L is local.
  }[tab:threats]
    \begin{tabular}{@{}lp{0.72\linewidth}@{}}
      \toprule
      Threat & Capability \\
      \midrule
      \threat{R1} & Guess the password online or replay a password leaked
        from another service (credential stuffing) \\
      \threat{R2} & Read the service's stored credentials \\
      \threat{R3} & Trick the user into typing credentials into a fake
        site, to use later (phishing) \\
      \threat{R4} & Relay the user's login to the real site as it happens
        (real-time phishing, adversary in the middle) \\
      \threat{R5} & Take over the channel of the second factor (the phone
        number, the email account) \\
      \threat{R6} & Talk or tire the user into approving a login \\
      \threat{R7} & Steal the session after a successful login \\
      \threat{R8} & Use account recovery or a weaker fallback login \\
      \midrule
      \threat{L1} & Hold the user's lost or stolen device \\
      \threat{L2} & Use a device shared with the user, or watch the user
        log in \\
      \threat{L3} & Run malware on the user's device \\
      \bottomrule
    \end{tabular}
  \end{sidecaption}
\end{table}

\subsection{Necessary and sufficient}\label{sec:logic}

Multi-factor authentication is \emph{sufficient} for security against a
set of threats if every multi-factor scheme resists every threat in the
set.
It is \emph{necessary} if no scheme with a single factor resists them.
Both are universal claims, so they are refuted by counterexamples of
different kinds.
One documented attack that succeeds against a multi-factor scheme refutes
sufficiency for that scheme and that threat.
Necessity is refuted by a single-factor scheme that resists the threats
at least as well as multi-factor schemes do, which takes a comparison and
not only an attack.

Security in practice is a matter of degree, so we also ask the graded
versions: how much compromise remains with multi-factor authentication,
and how close single-factor schemes come.
The answers may differ between the remote and the local threats, and
between the cases of \cref{tab:cases}; we report them separately.

\subsection{Research questions and hypotheses}\label{sec:questions}

The first question is a precondition for the others: unless multi-factor
authentication has an effect, there is nothing to attribute.
\begin{restatable}{question}{rqeffect}\label{rq:effect}
  Does multi-factor authentication reduce account compromise, by how
  much, and for which kinds of second factor?
\end{restatable}
\begin{restatable}{hypothesis}{heffect}\label{h:effect}
  Multi-factor authentication reduces account compromise, and the size
  of the reduction differs widely between kinds of second factor.
\end{restatable}

The second question separates the two explanations of that effect.
\begin{restatable}{question}{rqattribution}\label{rq:attribution}
  Is the effect due to combining factors or to the strength of one of
  them?
  What does the password contribute once a cryptographic key is present,
  and what does a second factor contribute beside a strong, unique
  password?
\end{restatable}
\begin{restatable}{hypothesis}{hattribution}\label{h:attribution}
  Against the remote threats the effect comes from the strongest factor:
  beside a cryptographic key bound to the service, the password adds
  nothing.
  Against the local threats a second, locally verified factor does add
  protection.
\end{restatable}

The last two questions are the ones in the title.
\begin{restatable}{question}{rqsufficient}\label{rq:sufficient}
  Is multi-factor authentication sufficient: are accounts protected by
  it secure against the threats of \cref{tab:threats}?
\end{restatable}
\begin{restatable}{hypothesis}{hsufficient}\label{h:sufficient}
  Multi-factor authentication is not sufficient: for each of the cases
  \case{A1}--\case{A3} and \case{B} at least one threat of
  \cref{tab:threats} succeeds.
\end{restatable}
\begin{restatable}{question}{rqnecessary}\label{rq:necessary}
  Is multi-factor authentication necessary: can comparable security be
  reached with a single factor, or by other means?
\end{restatable}
\begin{restatable}{hypothesis}{hnecessary}\label{h:necessary}
  Multi-factor authentication is not necessary against the remote
  threats: a single cryptographic key (case \case{C}) resists them as
  well as the multi-factor cases do.
\end{restatable}

The hypotheses are predictions, written down so that the search can
contradict them.
They are our expectation before reading; the searches are designed to
find evidence on both sides of each (\cref{sec:queries}).

\subsection{Contributions}

The report contributes the following; each contribution is held to
account by a research question.
\begin{itemize}
  \item An account of how large the measured effect of multi-factor
    authentication is and how well established (\cref{rq:effect}).
  \item A comparison, threat by threat, of what each factor contributes
    in the cases of \cref{tab:cases}
    (\cref{rq:attribution}).
  \item Answers, with their conditions, to whether multi-factor
    authentication is sufficient (\cref{rq:sufficient}) and necessary
    (\cref{rq:necessary}).
\end{itemize}


\section{Method}\label{sec:method}

\subsection{Overview}\label{sec:overview}

Each question is answered from the literature, by a search in both
directions: for work that supports the hypothesis and for work that
contradicts or narrows it.

\rqeffect*
Answered from studies that measure compromise with and without
multi-factor authentication, by testing \cref{h:effect}:
\heffect*

\rqattribution*
Answered from comparisons between kinds of factor, from evaluations of
passwordless and single-factor deployments, and from analyses of how
factors compose; summarised threat by threat for the cases of
\cref{tab:cases}, by testing \cref{h:attribution}:
\hattribution*

\rqsufficient*
Answered from documented attacks on accounts protected by multi-factor
authentication and from analyses that claim such schemes secure, by
testing \cref{h:sufficient}:
\hsufficient*

\rqnecessary*
Answered from evaluations of alternatives (single cryptographic keys,
risk-based authentication, password managers) and from what the
standards require at which level and why, by testing \cref{h:necessary}:
\hnecessary*

\subsection{Type of review}\label{sec:reviewtype}

This is a rapid review with a scoping aim: it maps what kinds of evidence
exist for four related questions and reports what that evidence says.

The aim is that of a scoping review, whose purpose is
\textcquote{Munn2018ScopingOrSystematic}{to identify and map the
available evidence}.
Such a review suits a topic where it is
\textcquote{Munn2018ScopingOrSystematic}{still unclear what other, more
specific questions can be posed and valuably addressed by a more precise
systematic review}.
That is the situation here: the title asks one question, but answering it
takes four, and the evidence for them is of different kinds.

The conduct is that of a rapid review.
There is no formal definition of one; a working definition from health
research, where the method comes from, is
\textcquote{Tricco2015RapidReviews}{a form of knowledge synthesis in
which components of the systematic review process are simplified or
omitted to produce information in a timely manner}.
The same notion is used in software engineering
research \parencite{Cartaxo2020RapidReviewsSE}.

The search is nevertheless systematic in the sense that it is planned,
documented and repeatable, which is what the course reading on rigour
asks of any review:
\textcquote[p.~163]{Bandara2015Rigor}{Defining a search strategy
significantly contributes to methodical rigor; it provides transparency,
clarity, and replicability.}

It is not a systematic review.
Three shortcuts separate it from one, and each limits what may be
concluded.
\begin{description}
  \item[Capped search] At most 40 records are taken per database and
    query, so the search samples the most relevant records and does not
    exhaust the literature.
    Shortcuts in the search lose information: when 2512 systematic
    reviews of clinical trials were repeated with rapid methods, the
    statistical significance of the result changed in 6.5 to 38.6 per
    cent of the meta-analyses, depending on the shortcut, although no
    bias in a particular direction was
    found \parencite{Marshall2019RapidVsSystematic}.
  \item[One screener] Records are screened once, by a language model
    whose doubtful decisions are checked by reading, and not by two
    independent human reviewers.
    Among human reviewers a single screener misses relevant studies, a
    median of 5 per cent and at worst 58 per cent in the evaluations
    reviewed by \textcite{Waffenschmidt2019SingleScreening}.
    We know of no corresponding figure for the language model used here.
  \item[Graded, not appraised] Each included study is graded by its
    design (\cref{sec:grading}); no instrument for appraising the quality
    of individual studies is applied.
\end{description}
A narrative review, reading what one already knows of, would have been
quicker still.
It was rejected because it cannot show that contradicting work was
looked for, which is the point of a question about necessity and
sufficiency.
A full systematic review was rejected for a first answer because the
questions are broad and the evidence is of mixed kinds; it is the natural
next step for the questions this review finds to be open
(\cref{sec:future}).

\subsection{Databases and tool}\label{sec:sources}

The searches are run with the tool \texttt{scholar}, which sends one
query to several databases and records the query, the date and the
number of records returned by each.
The databases are OpenAlex, Web of Science, Scopus, IEEE Xplore and
arXiv.
Two further databases the tool can reach, Semantic Scholar and DBLP, did
not answer on the day of the search and are therefore not searched.
Every query is run on every one of the five databases, translated into
the field syntax of the databases that need it but with the same
keywords.

Standards, guidance from public authorities and measurement reports from
service providers are often not indexed in these databases.
They are looked for with a web search, whose queries are recorded, and
then retrieved from the issuing organisation.

\subsection{Queries}\label{sec:queries}

All queries share the block of terms
\begin{quote}
  $M$ = \q{multi-factor authentication} \kw{OR}
  \q{multifactor authentication} \kw{OR}
  \q{two-factor authentication} \kw{OR} \q{2FA} \kw{OR}
  \q{two-step verification}.
\end{quote}
The bare abbreviation MFA is left out on purpose: it also stands for
material flow analysis, and a search with it returns studies of waste
recycling.
The block for a single strong factor is
\begin{quote}
  $K$ = \kw{passwordless} \kw{OR} \kw{passkey} \kw{OR} \kw{FIDO2}
  \kw{OR} \kw{WebAuthn} \kw{OR} \q{security key} \kw{OR}
  \q{single-factor}.
\end{quote}
\Cref{tab:planned} gives the planned queries.
For each question there is at least one query worded to find support for
the hypothesis and one worded to find the opposite, and both use the
same concept terms.
A query may be reworded after a first run if it returns mostly records
on another topic; every such change is reported in \cref{sec:deviations}
together with the query it replaced.
\begin{table}
  \begin{sidecaption}{%
    Planned queries, written before the search.
    $M$ and $K$ are the term blocks defined in \cref{sec:queries}.
    S marks a query worded to find support for the hypothesis of the
    question, C one worded to find evidence against it.
    The numbers of records found are given per database in
    \cref{sec:hitlists}.
  }[tab:planned]
    \footnotesize
    \begin{tabular}{@{}llp{0.68\linewidth}@{}}
      \toprule
      Question & Query & Keywords \\
      \midrule
      \Cref{rq:effect}
      & S1 & $M$ \kw{AND} (\q{account compromise} \kw{OR} \q{account
        takeover} \kw{OR} \kw{hijacking}) \kw{AND} (\kw{effectiveness}
        \kw{OR} \kw{empirical} \kw{OR} \kw{measurement}) \\
      & S2 & $M$ \kw{AND} (\kw{adoption} \kw{OR} \kw{deployment})
        \kw{AND} (\q{risk reduction} \kw{OR} \q{compromise rate} \kw{OR}
        \kw{incident}) \\
      & C1 & $M$ \kw{AND} (\kw{ineffective} \kw{OR} \q{no effect} \kw{OR}
        \q{does not prevent} \kw{OR} \q{false sense of security} \kw{OR}
        \kw{limitations}) \\
      \midrule
      \Cref{rq:attribution}
      & S1 & ($K$) \kw{AND} \kw{password} \kw{AND} (\kw{replace} \kw{OR}
        \kw{comparison} \kw{OR} \q{security analysis} \kw{OR}
        \kw{evaluation}) \\
      & S2 & ($K$) \kw{AND} (\q{user verification} \kw{OR} \kw{PIN}
        \kw{OR} \kw{biometric}) \kw{AND} (\kw{stolen} \kw{OR} \kw{theft}
        \kw{OR} \q{local attack}) \\
      & C1 & $M$ \kw{AND} (\q{defence in depth} \kw{OR} \q{defense in
        depth} \kw{OR} \q{independent factors} \kw{OR} \q{factor
        independence} \kw{OR} \q{combination of factors}) \\
      & C2 & ($K$) \kw{AND} \kw{password} \kw{AND} (\q{still needed}
        \kw{OR} \kw{fallback} \kw{OR} \kw{weakness} \kw{OR}
        \kw{vulnerability}) \\
      \midrule
      \Cref{rq:sufficient}
      & S1 & $M$ \kw{AND} (\kw{bypass} \kw{OR} \kw{phishing} \kw{OR}
        \q{man-in-the-middle} \kw{OR} \q{adversary-in-the-middle}) \\
      & S2 & $M$ \kw{AND} (\q{SIM swap} \kw{OR} \q{push fatigue} \kw{OR}
        \q{prompt bombing} \kw{OR} \q{session hijacking} \kw{OR}
        \q{account recovery}) \\
      & C1 & $M$ \kw{AND} (\q{formal verification} \kw{OR} \q{provable
        security} \kw{OR} \q{security proof} \kw{OR} \q{phishing-resistant}) \\
      \midrule
      \Cref{rq:necessary}
      & S1 & ($K$ \kw{OR} \q{risk-based authentication} \kw{OR}
        \q{password manager}) \kw{AND} (\kw{alternative} \kw{OR}
        \q{instead of} \kw{OR} \kw{without}) \kw{AND} ($M$) \\
      & S2 & \q{risk-based authentication} \kw{AND} (\kw{evaluation}
        \kw{OR} \kw{effectiveness} \kw{OR} \kw{security}) \\
      & C1 & $M$ \kw{AND} (\kw{necessary} \kw{OR} \kw{mandatory} \kw{OR}
        \kw{mandate} \kw{OR} \kw{required}) \kw{AND} (\kw{security}
        \kw{OR} \kw{compromise}) \\
      & C2 & \q{password-only} \kw{OR} \q{single-factor authentication}
        \kw{AND} (\kw{insufficient} \kw{OR} \kw{inadequate} \kw{OR}
        \kw{compromise}) \\
      \bottomrule
    \end{tabular}
  \end{sidecaption}
\end{table}

\subsection{Inclusion and exclusion}\label{sec:criteria}

A record is included for a question if it bears on that question for the
authentication of human users to online services, and reports evidence
of one of the kinds in \cref{tab:grades}.
There is no limit on the year of publication; the year is reported,
because the schemes in use have changed.
Only records in English are included.
For reviews of clinical trials that restriction has been found to change
few conclusions \parencite[\abso]{NussbaumerStreit2020NonEnglish}; we
know of no such test for security research.

A record is excluded if it
\begin{itemize}
  \item proposes a new authentication scheme and offers only its own
    security argument for that scheme, without bearing on the general
    question;
  \item concerns authentication between machines or devices, with no
    human user;
  \item studies only the usability or adoption of multi-factor
    authentication and reports nothing about security;
  \item gives an opinion or a recommendation without evidence or a
    stated method.
\end{itemize}
Records excluded under the first three points are counted as adjacent
work and not as evidence on either side.

\subsection{Screening and reading}\label{sec:screening}

The whole review was carried out by a language model (Claude, through
Claude Code) under the direction of the author: it ran the searches,
screened the records, read the sources and drafted the text.
The steps were as follows.
\begin{enumerate}
  \item All records returned for a question are collected in one session
    of the search tool, and missing abstracts are added.
  \item Each record is classified from its title and abstract into
    exactly one of four categories: supports the hypothesis, contradicts
    or narrows it, adjacent work, or another topic.
    The classification is made by a language model against the question
    and the criteria of \cref{sec:criteria}, with a stated confidence.
  \item Every record classified as contradicting or narrowing, and every
    record classified with low confidence, is read again and reclassified
    where the first classification was wrong.
  \item The records that bear on the question are read in full text
    where it can be obtained.
    For each source cited, a passage that states the finding is
    recorded verbatim with its location.
    A source that could be read only as an abstract is marked as such
    where it is cited.
  \item A source whose existence or content cannot be checked against
    the publisher's record is dropped and listed as dropped.
\end{enumerate}
The number of records at each step is reported per question.
The provenance of each cited source (the query that found it, why it was
chosen, the passage relied on, and the outcome of the search for
contradicting work) is recorded with its entry in the source file of the
bibliography.

\subsection{Grading the evidence}\label{sec:grading}

Sources are graded by the kind of evidence they give, in
\cref{tab:grades}.
The grade says what a source can show, not how good the study is.
The grades are not a ranking for every purpose.
An attack demonstration is weak evidence of how often accounts are
compromised but conclusive evidence that a scheme can be bypassed, which
is what \cref{rq:sufficient} asks.
A measurement of real accounts is the strongest evidence for
\cref{rq:effect} but, being observational, cannot by itself say why the
effect arises, which is what \cref{rq:attribution} asks.
\begin{table}
  \begin{sidecaption}{%
    Kinds of evidence.
    The last column says what a source of that kind can establish.
  }[tab:grades]
    \footnotesize
    \begin{tabular}{@{}llp{0.42\linewidth}@{}}
      \toprule
      Grade & Kind of evidence & Can establish \\
      \midrule
      E1 & Controlled experiment & That a difference is caused by the
        scheme, in the setting studied \\
      E2 & Measurement of real accounts or incidents & How often
        compromise occurs with and without a scheme; not why \\
      E3 & Attack, demonstrated or observed in use & That a scheme can be
        bypassed; not how often \\
      E4 & Formal or analytical security analysis & What follows from a
        stated model \\
      E5 & Standard or expert consensus & What is required or
        recommended, and the reasons given \\
      E6 & Provider's report without disclosed method & That the provider
        makes the claim \\
      \bottomrule
    \end{tabular}
  \end{sidecaption}
\end{table}

\subsection{Synthesis}\label{sec:synthesis}

For \cref{rq:effect} the reported effect sizes are tabulated with their
populations, the kind of second factor and the grade of evidence.
For \cref{rq:attribution,rq:sufficient,rq:necessary} the findings are
entered in one table of the cases of \cref{tab:cases} against the threats
of \cref{tab:threats}.
A cell says whether the case resists the threat, and cites the source;
a cell for which the search found no evidence is left open and reported
as such.
The answers to the questions are read off that table: sufficiency fails
for a case if any cell in its row is a successful attack, and necessity
fails for a set of threats if a single-factor row resists all of them.

\subsection{Deviations from the protocol}\label{sec:deviations}

Everything above this subsection was written before the search, with
three exceptions, all of them added sources.
The two sources for the premise in the first paragraph of the
introduction were found by the search.
The sources on types of review in \cref{sec:reviewtype} and the remark
on the language restriction in \cref{sec:criteria} come from a fifth,
small search on review methods that the protocol had not planned
(\cref{sec:hits-method}).
The search itself departed from the protocol as follows.
\begin{description}
  \item[Queries added] Six queries were added, for three of the
    questions, because the planned ones returned few records on the cases
    with a cryptographic key or mostly proposals of new schemes.
    They use the same term blocks and are marked as added in the tables
    of \cref{sec:hitlists}; the planned queries were kept.
  \item[Searches that failed] IEEE Xplore allows a fixed number of
    requests a day.
    The planned queries were all answered, but four of the added ones
    were not, and are reported as failed.
  \item[Screening model] The search tool's default language model was
    not available.
    Claude Sonnet 5.5 classified the records for three questions and for
    the search on review methods, Claude Opus 4.8 those for
    \cref{rq:effect}.
    For \cref{rq:effect} the model filed every measurement study under
    \enquote{contradicts or narrows}, apparently because the studies also
    report differences between factors; 52 records were reclassified by
    reading.
    For the other questions 14, 15 and 9 records were reclassified.
  \item[Records of low confidence] The protocol promised that every
    record classified with low confidence would be read again.
    That was done for the records classified as bearing on a question.
    Records classified as adjacent with low confidence were checked on
    their titles, and only some on their abstracts.
  \item[A second check of quotations] The four questions were searched
    and read by four separate runs of the model.
    A fifth run then checked every quoted passage again against the full
    text, or against the abstract for sources read only as abstracts.
    The check found one misquoted hyphen and led to three sources being
    dropped (\cref{sec:dropped}).
  \item[Sources read only as abstracts] Of the 64 sources cited, 21 could
    be read only at the level of the abstract, mostly because the
    publisher's site refused automated retrieval.
    They are marked \enquote{\abso} where they are cited.
  \item[Synthesis tables] The cells of the tables of cases against
    threats hold a symbol and not a citation; the sources are cited in
    the text, threat by threat.
    No query had terms for threats \threat{R2} and \threat{L2}, and the
    search found no evidence on them, so those columns are open.
  \item[Grades] Two runs graded one source differently; it was set to
    the grade that fits its method (an evaluation by analysis, E4).
\end{description}


\section{Results}\label{sec:results}

\subsection{What the search returned}\label{sec:flow}

\Cref{tab:flow} gives the number of records at each step.
Most records returned by a query that names multi-factor authentication
are proposals of new schemes, which the criteria of \cref{sec:criteria}
count as adjacent.
Between 28 and 45 records per question bear on it.
\begin{table}
  \begin{sidecaption}{%
    Records at each step, per question.
    \enquote{Supporting} and \enquote{contradicting} are relative to the
    hypothesis of the question.
    Distinct records are counted after merging versions of the same
    work.
    Sources cited include standards and reports found outside the
    databases; a source cited for two questions is counted under both.
  }[tab:flow]
    \footnotesize
    \begin{tabular}{@{}lrrrr@{}}
      \toprule
      & \multicolumn{4}{c}{Question} \\
      \cmidrule(l){2-5}
      Step & 1 & 2 & 3 & 4 \\
      \midrule
      Records returned by the databases & 627 & 522 & 780 & 772 \\
      Distinct records & 387 & 289 & 508 & 497 \\
      \quad classified as adjacent & 290 & 208 & 342 & 406 \\
      \quad classified as another topic & 69 & 48 & 123 & 46 \\
      \quad bearing on the question & 28 & 33 & 43 & 45 \\
      \qquad cited & 9 & 15 & 19 & 18 \\
      \qquad supporting, not cited & 0 & 7 & 16 & 9 \\
      \qquad contradicting or narrowing, not cited & 17 & 10 & 7 & 18 \\
      \qquad excluded after reading & 2 & 1 & 1 & 0 \\
      \midrule
      Sources cited, with those outside the databases & 11 & 19 & 25 & 24 \\
      \quad of which read only as abstract & 5 & 8 & 4 & 8 \\
      \bottomrule
    \end{tabular}
  \end{sidecaption}
\end{table}

\subsection{Effect}\label{sec:res-effect}

Five studies measure compromise of real accounts with and without a
second factor (grade E2), and two providers state a figure without a
method (E6).
\Cref{tab:effects} lists them.
All seven point the same way: less compromise with a second factor.
\begin{table}
  \begin{sidecaption}{%
    Reported effects of a second factor on account compromise.
    \enquote{Abs.} marks a source read only as an abstract.
    The figures measure different things and cannot be pooled.
  }[tab:effects]
    \footnotesize
    \begin{tabular}{@{}p{0.19\linewidth}p{0.25\linewidth}p{0.30\linewidth}l@{}}
      \toprule
      Source & Setting & Reported effect & Grade \\
      \midrule
      Meyer et al.\ 2023
        & Microsoft business accounts under review for suspicious
          activity
        & 99.22\% lower risk of compromise; 98.56\% when the password
          had leaked
        & E2 \\
      Doerfler et al.\ 2019
        & Google accounts, 350\,000 hijacking attempts
        & Device-based challenges blocked over 94\% of attempts from
          phishing and 100\% of automated ones
        & E2 \\
      Kamar et al.\ 2026 (abs.)
        & Three Canadian banks
        & Mandatory: fewer compromised accounts; optional: no
          significant change
        & E2 \\
      Kawase et al.\ 2019 (abs.)
        & One online marketplace
        & 43\% fewer account takeovers
        & E2 \\
      Egenti and Ojo 2026 (abs.)
        & Online banking in Nigeria
        & Fraud cases down 10 to 15\%
        & E2 \\
      Google 2022
        & Google accounts enrolled automatically
        & 50\% fewer compromised accounts
        & E6 \\
      Microsoft 2019
        & Microsoft cloud services
        & Blocks over 99.9\% of account compromise attacks
        & E6 \\
      \bottomrule
    \end{tabular}
  \end{sidecaption}
\end{table}

The two largest studies are by the providers' own staff.
\Textcite{Meyer2023MFAEffect} estimate, for Microsoft's business
customers, that
\textcquote[abstract]{Meyer2023MFAEffect}{MFA reduces the risk of
compromise by 99.22\% across the entire population and by 98.56\% in
cases of leaked credentials}.
The sample is accounts that were reviewed because of suspicious
activity, and whether an account had a second factor was decided by its
organisation, not assigned at random.
\Textcite{Doerfler2019LoginChallenges} study challenges that Google
shows when a login looks risky, which is not the same as a second factor
required at every login:
\textcquote[abstract]{Doerfler2019LoginChallenges}{Device-based
challenges provide the best protection, blocking over 94\% of hijacking
attempts rooted in phishing and 100\% of automated hijacking attempts}.
There is no group of accounts without challenges to compare with.

The studies from outside the large providers report smaller effects.
One marketplace reports 43 per cent fewer account takeovers after
introducing a second factor \parencite[\abso]{Kawase2019AccountTakeover}.
At three Canadian banks,
\textcquote[abstract; \abso]{Kamar2026TwoFactorBanks}{mandatory 2FA
significantly reduced the number of compromised bank accounts, whereas
optional 2FA did not}.
A study of Nigerian online banking reports that fraud cases fell by 10
to 15 per cent \parencite[\abso]{Egenti2026Nigeria}; its figures come
partly from a prototype and we give it little weight.
The two provider statements, a 50 per cent decrease at Google
\parencite{Google2022SaferInternet} and over 99.9 per cent of attacks
blocked at Microsoft \parencite{Maynes2019OneSimpleAction}, give no
method.

Two studies separate kinds of second factor.
At Microsoft, a code by text message was
\textcquote[sec.~5]{Meyer2023MFAEffect}{40.8\% less effective than
Microsoft Authenticator}, an app, although the failure rates were low
for both (1.66 against about 1 per cent).
At Google, codes by text message were
\textcquote[sec.~1]{Doerfler2019LoginChallenges}{preventing only 96\% of
attacks rooted in phishing and 76\% of targeted attacks}, against 99 and
90 per cent for prompts on the user's device, and
\textcquote[sec.~4.1]{Doerfler2019LoginChallenges}{security keys provide
the best protection, blocking 100\% of attacks in our dataset}.
Challenges that ask for something the user knows, such as a secondary
email address, stopped only 10 per cent of the attempts from phishing.

Several sources narrow the result.
The null result for optional two-factor authentication above is the only
measurement in which a second factor made no difference.
Services that take over accounts for hire were found able
\textcquote[abstract; \abso]{Mirian2019HackForHire}{to bypass SMS
two-factor authentication}.
Flaws in how services implement the second step let an attacker skip it
at three providers \parencite{Wang2024SimpleNotSecure}, and a criminal
service sells stolen user profiles with which buyers
\textcquote[abstract]{Campobasso2020IMPaaS}{effectively bypass
multi-factor authentication mechanisms}.
These show that a bypass is possible, not how often it happens.

The search worded to find the opposite, a study in which multi-factor
authentication made no difference or made things worse, returned none
apart from the optional case above.
That is the outcome of a capped search and not proof that no such study
exists.

In sum, \cref{h:effect} holds in its direction.
Its second half, that the size differs widely between kinds of second
factor, holds only in part: within one provider the kinds differ by a
few percentage points of the attempts, while the spread between studies,
from 43 to over 99 per cent, mixes the kind of factor with the setting,
the outcome measured and who did the measuring.
No study found measures compromise for a cryptographic key without a
password (cases \case{B} and \case{C}) against a password with a second
factor.

\subsection{How the standards classify the cases}\label{sec:res-standards}

The working definitions of \cref{sec:terms} agree with the standards.
The guideline of the United States standards institute defines
multi-factor authentication by kinds of factor, and at its second
assurance level
\textcquote[sec.~2.2.1]{NIST2025SP80063B}{authentication SHALL use
either a multi-factor authenticator or a combination of two separate
authentication factors}.

All three standards and assessments read classify a key with user
verification (case \case{B}) as multi-factor and a key without it (case
\case{C}) as single-factor.
In the guideline, a key that is unlocked by a PIN or biometric is a
multi-factor authenticator, and if the key reports that it was not
unlocked in that way the authentication
\textcquote[sec.~3.1.7.2]{NIST2025SP80063B}{SHALL be treated as
single-factor}.
The web standard for such keys says the same:
\textcquote[sec.~6.2.3]{W3C2026WebAuthn3}{an authenticator that supports
user verification is multi-factor capable}.
The United Kingdom's cyber security centre agrees and answers the
objection that key and verification sit in one device: factors are kinds
of evidence, not devices \parencite{NCSC2026CredentialsPaper}.

Passwordless authentication with a passkey is therefore, in the
standards' terms, usually multi-factor authentication.
What the standards attach to the key itself, whether or not it is
unlocked by a second factor, is resistance to phishing:
\textcquote[sec.~3.2.5]{NIST2025SP80063B}{Phishing resistance requires
single- or multi-factor cryptographic authentication.}

\subsection{Cases against threats}\label{sec:res-table}

\Cref{tab:remote,tab:local} summarise what the sources say about each
case of \cref{tab:cases} against each threat of \cref{tab:threats}.
The rest of this subsection gives the sources, threat by threat.
The one source that assesses most cases against most threats in one
place is a technical paper by the United Kingdom's cyber security centre
\parencite{NCSC2026CredentialsPaper}.
It is an assessment by an authority (E5): it does not say how each
verdict was reached, and its authors recommend passkeys.
Where it is the only source for a cell we say so.
\begin{table}
  \begin{sidecaption}{%
    Cases against remote threats.
    \res{} resists, \fail{} fails, \partly{} depends on conditions given
    in the text, \na{} the threat does not apply, empty: no evidence
    found.
    The row for case \case{C} rests on the standards' classification, not
    on tests.
  }[tab:remote]
    \footnotesize
    \begin{tabular}{@{}lcccccccc@{}}
      \toprule
      Case & \threat{R1} & \threat{R2} & \threat{R3} & \threat{R4}
        & \threat{R5} & \threat{R6} & \threat{R7} & \threat{R8} \\
      \midrule
      \case{P}  & \fail & & \fail & \fail & \na & & \fail & \\
      \case{P+} & \res  & & \fail & \fail & \na & & \fail & \\
      \case{A1} & \res  & & \res  & \fail & \fail & & \fail & \partly \\
      \case{A2} & \res  & & \res  & \fail & & \fail & \fail & \partly \\
      \case{A3} & \res  & & \res  & \res  & \na & & \fail & \partly \\
      \case{B}  & \res  & & \res  & \res  & \na & \res & \fail & \partly \\
      \case{C}  &       & & \res  & \res  & \na & & & \partly \\
      \bottomrule
    \end{tabular}
  \end{sidecaption}
\end{table}
\begin{table}
  \begin{sidecaption}{%
    Cases against local threats; symbols as in \cref{tab:remote}.
    No source tests case \case{C}.
  }[tab:local]
    \footnotesize
    \begin{tabular}{@{}lccc@{}}
      \toprule
      Case & \threat{L1} & \threat{L2} & \threat{L3} \\
      \midrule
      \case{P}  & \res    & & \fail \\
      \case{P+} & \fail   & & \fail \\
      \case{A1} & \partly & & \partly \\
      \case{A2} & \fail   & & \fail \\
      \case{A3} & \partly & & \fail \\
      \case{B}  & \partly & & \fail \\
      \case{C}  &         & & \\
      \bottomrule
    \end{tabular}
  \end{sidecaption}
\end{table}

\paragraph{Guessing and reuse (\threat{R1})}
A password chosen and reused by the user falls to credential stuffing;
one analysed breach began with about 14\,000 accounts taken over in that
way \parencite[\abso]{Holthouse2025TwentyThreeAndMe}.
A random password unique to the service does not, and neither does any
case with a second factor or a key, in the assessment of
\textcite{NCSC2026CredentialsPaper}.
The measurements agree for the second factor: the effect at Microsoft
was 98.56 per cent for accounts whose password had already leaked
\parencite{Meyer2023MFAEffect}.

\paragraph{Phishing (\threat{R3}, \threat{R4})}
A password alone, however strong, can be phished.
A code or a push approval stops the use of a phished password later but
not a login relayed as it happens: the cyber security centre calls all
such second factors
\textcquote[summary]{NCSC2026CredentialsPaper}{inherently phishable}, and
the guideline states that
\textcquote[sec.~3.1.3]{NIST2025SP80063B}{Out-of-band authentication is
not phishing-resistant}.
A laboratory relay captured the session through codes by text message,
call, email and app and through push approvals, and its author judges
\textcquote[abstract]{Kikuchi2025Evilginx}{that 80\% of major web
services are vulnerable}; a small controlled campaign found the same for
simulated logins \parencite{Cuellar2026AiTM}.
A cryptographic key bound to the service resists both threats in every
source that examines it: a formal analysis that assumes the password
already lost \parencite{Jacomme2021FormalMFA}, a proof of the web
protocol's authentication property \parencite{Barbosa2021ProvableFIDO2},
and an analysis of attacks concluding that
\textcquote[abstract]{Berladskyy2026AttackVectors}{The claim that
passkeys are phishing-resistant largely holds true}.
Two providers report the same from use, without a method: at one,
attackers who had employees' passwords
\textcquote{Cloudflare2022Phishing}{could not get past the hard key
requirement}, and at the other
\textcquote{Thomas2019Hygiene}{zero users that exclusively use security
keys fell victim to targeted phishing}.
One demonstrated attack defeats keys too, when the service has a flaw
that lets the attacker run a script in its pages
\parencite[\abso]{Catalano2025BitM}.

\paragraph{Taking over the channel (\threat{R5})}
A public review of a series of real intrusions found that
\textcquote[p.~5]{CSRB2023Lapsus}{attackers gained initial access to
targeted organizations through Subscriber Identity Module (SIM) swapping
attacks, which allowed them to intercept one-time passcodes and push
notifications sent via SMS, effectively defeating this widely used MFA
control}.
The threat applies only where the second factor travels over a channel
that can be taken over.

\paragraph{Approval by the user (\threat{R6})}
With the password known, an attacker who triggers a push approval was
let in by the user in about 68 per cent of the sessions of one
experiment \parencite[\abso]{Jubur2021PushBypass}.
The United States cybersecurity agency lists such
\enquote{push bombing} among the weaknesses of some forms of
multi-factor authentication \parencite{CISA2022PhishingResistant}.
The cyber security centre rates keys as not exposed to it.
We found no report of the frequency of this attack in use.

\paragraph{Session theft (\threat{R7})}
What is stolen here is the session that follows a successful login, so
no authenticator prevents it; the cyber security centre notes that
common malware copies session cookies as well as passwords, and that a
relayed login leaves the session with the attacker
\parencite{NCSC2026CredentialsPaper}.
The criminal service above trades in such stolen sessions
\parencite{Campobasso2020IMPaaS}, and a review describes cookie theft as
a service \parencite[\abso]{Pazmino2025Cookie}.

\paragraph{Recovery and fallback (\threat{R8})}
An account is as strong as its weakest way in.
On popular sites, a second factor was in some cases switched off on
request when the user said it was lost \parencite{Gerlitz2023Recovery},
although in a test of ten services a password reset alone did not get
past the second factor \parencite{Buttner2025Inconsistencies}.
For keys, a login can be pushed down to a weaker method where one is
still enrolled; a vendor demonstrated this and adds that
\textcquote{Miron2025Downgrade}{there is currently no evidence for the
usage of this technique by threat actors in the wild}.
The methods for recovering an account after losing a key
\textcquote[abstract]{Kunke2021Recovery}{have many drawbacks. Some even
rely on passwords, taking passwordless authentication ad absurdum.}
Recovery remains
\textcquote[conclusion]{Buttner2024Recovery}{a relevant entry point for
account takeover attacks}, and a study of account access in use reports
\textcquote[abstract; \abso]{Wang2025TwoFAIllusion}{overlooked
weaknesses of current deployments}.

\paragraph{A lost or stolen device (\threat{L1})}
Here the cyber security centre's assessment is the main source.
It rates passwords kept in a password manager, codes from an app or by
text message, push approvals and passkeys on a phone as exposed when the
device is stolen and its PIN is weak, and a separate hardware token as
not exposed, whether the token produces codes or holds a key unlocked by
a PIN or password.
A password that the user remembers is the exception among the secrets:
\textcquote[note to table~19]{NCSC2026CredentialsPaper}{theft of a
device even with the PIN won't reveal that password}.
For a key bound to the device, what stands between a thief and the
account is user verification:
\textcquote[sec.~3.1]{Buttner2025DeviceBound}{it will be difficult to
steal a user's device-bound passkey, provided it is protected by user
verification mechanisms}.
That protection has been broken in particular products: extracting the
PIN was shown to be straightforward in practice
\parencite[\abso]{Barbosa2023RogueKey},
the protocol for the PIN was found to have flaws
\parencite{Guan2022FormalFIDO2}, and of the Windows configurations
tested
\textcquote[abstract; \abso]{Lee2025WindowsHello}{only systems with
Enhanced Sign-in Security (ESS) prevented our attacks}.

\paragraph{Malware on the device (\threat{L3})}
Every case fails against malware in at least one source.
A formal analysis finds hardware tokens secure against phishing but
concludes that
\textcquote[sec.~5.5]{Jacomme2021FormalMFA}{malware attacks are still
possible}.
Seven attacks through a malicious browser extension were demonstrated on
ten sites that use keys; the authors add that their analysis of existing
extensions
\textcquote[abstract]{Yadav2024LocalAttacks}{finds no evidence of the
attacks in the wild}.
A Trojan took over passkey logins in every trial of one laboratory
study \parencite{Mitra2026WoodenHorse}, and with a deceptive overlay
\textcquote[abstract; \abso]{Mahdad2024FIDOLA}{users approved
approximately 95.55\% of cross-service attacks}.
The cyber security centre is milder: it rates passkeys as
\enquote{situation dependent} and a hardware token with a PIN as not
exposed to malware, which the demonstrations above contradict.
For codes it separates those produced on the infected device, which
fail, from those of a separate token, which do not.
Three further analyses agree that the device is the weak point: one
concludes that
\textcquote[abstract]{Mitra2026TrustBoundary}{the weak link in a FIDO2
deployment is rarely the crypto layer}, one that passwords and keys are
both exposed to malicious scripts and extensions
\parencite[\abso]{Gautam2025Secret}, and one that second factors
believed secure against malware on the terminal are not
\parencite[\abso]{Mahdad2023SoK2FA}.

\subsection{Attribution}\label{sec:res-attribution}

No study found isolates the contribution of one factor in measured
compromise.
The answer to \cref{rq:attribution} therefore rests on
\cref{tab:remote,tab:local}, that is on assessment, analysis and attack
demonstrations, read as three comparisons.

\paragraph{A second factor beside a strong password}
Compare the rows \case{P+} and \case{A1}--\case{A3}.
Against guessing and reuse the second factor adds nothing, because the
strong password already resists.
Against a phished password used later it adds protection.
Against a relayed login a code or push adds nothing and a key adds
everything.
A code by text message and a push each bring a threat of their own
(\threat{R5}, \threat{R6}).
Against the local threats a second factor helps when it is not on the
device that holds the password.
What a second factor adds thus depends on which factor it is, and for
the remote threats a weak one adds little to a strong password.
The large measured effects of \cref{sec:res-effect} do not contradict
this: they compare with passwords as users choose them, and the one
figure for passwords already known to the attacker, 98.56 per cent, is
the work of the second factor alone.

\paragraph{A password beside a key}
Compare the rows \case{A3}, \case{B} and \case{C}.
Against the remote threats no source shows the password adding
anything.
The formal analysis of hardware tokens takes as its
\textcquote[sec.~1]{Jacomme2021FormalMFA}{basic assumption} that the
password is already compromised, and finds the token secure against
phishing all the same; the attackers at Cloudflare held the passwords
and were stopped by the key.
The review board that examined the intrusions of \threat{R5} draws the
same conclusion and recommends
\textcquote[p.~6]{CSRB2023Lapsus}{passwordless solutions}.
Against the local threats the password is not idle.
It is the one secret that a stolen device does not hold, and one study
of local attacks found
\textcquote[sec.~I]{Yadav2024LocalAttacks}{two attacks (3 and 4) where
passwordless authentication presents more risk than passwords alone in
the presence of a compromised extension}.

\paragraph{User verification beside a key}
Compare the rows \case{B} and \case{C}.
The standards distinguish them, but no source tests a key alone.
That user verification is what protects a stolen key is stated by
analysis and, for weak PINs and particular products, contradicted by
the attacks under \threat{L1} above.

\Cref{h:attribution} is thus supported in its first half and corrected
in its second.
Against the remote threats the protection is that of the strongest
factor, and a password beside a key bound to the service adds nothing
that any source could show; this rests on analysis, assessment and two
providers' reports, not on measurement.
Against the local threats a second factor does add protection, as
predicted, but only if it is not kept on the same device, and the
factor that most plainly meets that condition is a remembered password.

\subsection{Sufficiency}\label{sec:res-sufficient}

By the rule of \cref{sec:synthesis}, multi-factor authentication is not
sufficient for a case if any cell in its row is a successful attack.
Every multi-factor row of \cref{tab:remote,tab:local} has one.

For a code or a push beside a password (cases \case{A1} and \case{A2})
the attacks are documented in use.
The review board cited above concludes from real intrusions that
\textcquote[p.~5]{CSRB2023Lapsus}{the multi-factor authentication (MFA)
implementations used broadly in the digital ecosystem today are not
sufficient for most organizations or consumers}.

For the cases with a key (\case{A3} and \case{B}) the evidence is of
another kind.
No source reports an attack that defeats the key itself from a
distance, and none reports an account protected by a key being taken
over in use.
The attacks that succeed go around the key: through a weaker method
left enrolled, through recovery, through the session after login, or
through malware on the user's device.
All are demonstrations in a laboratory or by a vendor, several of them
not yet peer reviewed, and two of their authors state that they found
no sign of use.
Formal results cover the protocol and not the account: the
authentication of the web protocol is proved, the protocol between
client and token
\textcquote[abstract]{Barbosa2021ProvableFIDO2}{can only be proved
secure in a weak sense}, and another analysis finds that the scheme
\textcquote[abstract]{Guan2022FormalFIDO2}{fails to achieve some strong
authentication properties}.
A later analysis of the current protocol versions finds that their
resistance to downgrade
\textcquote[abstract; \abso]{Bindel2023ProvableFIDO2}{can be improved}.
One demonstrated attack on the protocol between client and token deletes
the user's keys without logging in as the user
\parencite{Casagrande2025CTRAPS}.
We found no formal result that covers recovery, sessions and malware
for a whole account.
A proof is in any case weaker evidence here than it looks: of the many
multi-factor schemes published with a proof,
\textcquote[abstract; \abso]{Wang2023Proofs}{most of them have been
shortly found to be insecure}, and a recent analysis found
vulnerabilities in ten more \parencite{Ang2025Excavating}.
No source separates a key with a password from a key with user
verification, so the two rows rest on the same evidence.

\Cref{h:sufficient} is supported: no multi-factor case is sufficient
against all the threats of \cref{tab:threats}.
The strength of that conclusion differs.
It is strong for codes and push approvals, where the failures are
observed.
For keys it rests on demonstrated possibility, and it holds because the
threat model includes what happens around the authenticator.
Against phishing alone, the evidence found gives no reason to call a
key insufficient.

\subsection{Necessity}\label{sec:res-necessary}

By the rule of \cref{sec:synthesis}, multi-factor authentication is not
necessary against a set of threats if a single-factor row resists all
of them.
Three single-factor alternatives were examined.

\paragraph{A key alone}
The standards treat a key without user verification as a legitimate
single-factor authenticator.
The guideline's lowest level
\textcquote[sec.~2.1]{NIST2025SP80063B}{requires either single-factor or
multi-factor authentication}, and it attributes resistance to phishing
to the cryptographic authenticator whether single- or multi-factor
(\cref{sec:res-standards}).
The scheme
\textcquote[sec.~1]{Kunke2021Recovery}{allows strong cryptographic
single factor authentication}, and one analysis judges that
\textcquote[abstract; \abso]{AlQubaisi2020Rush}{password-based
authentication is less secure than FIDO SF authentication}.
But the search found no source that tests a key alone against a
multi-factor case, for any threat.
The comparisons that exist, above all the cyber security centre's, are
for a key with user verification, which is multi-factor.

\paragraph{Risk-based authentication}
Here the service judges each login from its context and asks for more
only when the login looks unusual.
Replaying attacks on a record of 31.3 million login attempts at a real
service, the largest study concludes:
\textcquote[p.~6:28]{Wiefling2022PumpUp}{we consider RBA an effective
measure to protect users from the majority of attacks involving stolen
passwords}.
It has limits.
At large services
\textcquote[abstract]{Makowski2023RBA}{several test cases rarely trigger
the RBA system}, the criminal service of \cref{sec:res-effect} defeats
it with stolen profiles \parencite{Campobasso2020IMPaaS}, phishing sites
can copy the browser fingerprints it relies on
\parencite[\abso]{Lin2022PhishSheep}, and on mobile devices those
fingerprints are too alike to tell users apart
\parencite[\abso]{Spooren2015Fingerprinting}.
Nor is it an alternative with one factor: when it does ask for more,
\textcquote[conclusion]{Wiefling2019ReallyYou}{verification codes are
currently the unwritten standard}, so it rations a second factor
without replacing it, and the guideline rules that such signals do not
\textcquote[sec.~2]{NIST2025SP80063B}{substitute for an authentication
factor}.

\paragraph{A strong, unique password}
It resists guessing and reuse and nothing else in \cref{tab:remote}.
The cyber security centre advises it together with a second step where
passkeys are not available \parencite{NCSC2026PasskeysPage}.
Two web searches for a measured comparison of such passwords with
multi-factor authentication returned no study.

\paragraph{What is required}
The authorities require more than one factor at their upper levels.
At the guideline's second level
\textcquote[sec.~2.2]{NIST2025SP80063B}{Proof of possession and control
of two distinct authentication factors through the use of secure
authentication protocols is required}; the United States government
tells its agencies that for staff
\textcquote[p.~5]{OMB2022ZeroTrust}{phishing-resistant MFA is required}
and gives as reason
\textcquote[p.~7]{OMB2022ZeroTrust}{that enterprise users are among the
most valuable targets for phishing}.
An earlier white paper likewise argues for multi-factor authentication
of privileged users in place of passwords
\parencite[\abso]{NIST2016PIVPrivileged}.
The United Kingdom's centre now recommends
\textcquote{NCSC2026PasskeysBlog}{passkeys wherever a service supports
them, and two-step verification (2SV) where it does not}, and advises
services against asking for a second factor at every login
\parencite{NCSC2026Methods}.
A requirement shows what experts agree on (E5), not that fewer factors
would fail.
The reason these texts give is phishing, and the property they ask for
is one that the same guideline grants a single-factor key.

\Cref{h:necessary} is neither refuted nor established.
The evidence supports a weaker statement: against the remote threats it
is the binding of a key to the service that protects, not the number of
factors, so the standards' own reasoning does not make several factors
necessary there.
The stronger statement, that a key alone resists the remote threats as
well as the multi-factor cases, is untested.
Two further facts limit it.
No case in \cref{tab:remote}, with one factor or several, resists
session theft or weak recovery, so there is no secure single-factor row
by the rule above, and no secure multi-factor row either.
And a key alone protects a lost device with nothing: that follows from
what case \case{C} is and needs no source, which is also why none of the
sources bothers to test it.


\section{Discussion and limitations}\label{sec:discussion}

The four answers fit one picture.
Against attackers at a distance, authentication is about as strong as
its strongest factor.
A second factor helps a weak password a great deal, which is what the
measurements show, and helps a strong password only against the threats
that the second factor itself resists.
Once a factor is present that cannot be phished, the others add nothing
that the sources could show.
A second factor does its own work in one place, when the device is in
the wrong hands, and there it must be something the device does not
hold.
Several factors is therefore the wrong thing to count: what matters is
which threats each factor resists and whether the factors can be lost
together.

This picture is our reading of evidence of unequal strength, and the
following limits apply.
\begin{description}
  \item[Kinds of evidence] Only the first question is answered from
    measurements of real accounts, and those are mostly the providers'
    own, without random assignment.
    The other three rest on demonstrations, analysis and one authority's
    assessment.
    That assessment fills many cells of
    \cref{tab:remote,tab:local} alone, gives no method, and comes from a
    body that recommends passkeys.
  \item[What was not found] No measurement isolates one factor.
    No source tests a key alone.
    Nothing was found on stolen server credentials (\threat{R2}) or
    shared devices (\threat{L2}), for which the queries had no terms.
    Absence from a capped search is weak evidence of absence.
  \item[The search] At most 40 records were taken per query and
    database, and many queries reached that limit.
    Two databases, Semantic Scholar and DBLP, could not be searched, and
    IEEE Xplore answered only the planned queries.
    One database returns no abstracts, so some records were classified
    on their titles.
  \item[Reading] A third of the sources cited were read only as
    abstracts, among them the only null result for the first question
    and the experiments on push approval and deceptive overlays.
    An abstract states what the authors conclude and cannot be checked
    against their data.
  \item[Screening] One language model screened each question, and its
    decisions were checked by the same kind of model.
    For the first question it misfiled every supporting study.
    Whatever a model misses systematically, a second run of a model is
    likely to miss as well; no human screened the records.
  \item[Recency] Several of the sources on attacks against keys are
    preprints from 2025 and 2026.
    Their findings have not been reviewed, and products may have been
    fixed since.
  \item[Unmeasured dimensions] The review asks only about security.
    Usability, availability and the cost of losing a factor decide what
    can be deployed, and adoption decides the effect, as the difference
    between mandatory and optional use suggests.
\end{description}


\section{Conclusions}\label{sec:conclusions}

Multi-factor authentication is not sufficient for secure authentication
of users, and the literature does not show it to be necessary against
attackers at a distance; it is needed, in a limited sense, against the
loss of the device.
The answers to the four questions follow.

\rqeffect*
Yes.
Every measurement found shows less compromise with a second factor.
The size is not well established: the two providers' measurements are
above 94 per cent, one provider's own statement is 50 per cent,
independent studies report between 10 and 43 per cent, and optional use
showed no effect at two banks.
Codes by text message do worse than apps and prompts, and keys best.
\Cref{h:effect} holds in direction, and in part as to size.

\rqattribution*
Against remote threats, to the strength of one factor.
No source shows a password adding protection beside a cryptographic key
bound to the service, and a weak second factor adds little beside a
strong password.
Against local threats a second factor adds protection if it is not kept
on the same device, and a remembered password is such a factor.
\Cref{h:attribution} is supported in its first half and corrected in
its second.
None of this rests on measurement.

\rqsufficient*
No.
Codes and push approvals are defeated in use, by relayed logins, by
taking over the phone number and by tired users.
Keys have not been shown defeated from a distance, but accounts
protected by them can be entered through fallback methods, recovery,
stolen sessions and malware, as demonstrations show.
\Cref{h:sufficient} is supported, strongly for codes and push approvals
and on demonstrated possibility for keys.

\rqnecessary*
Not shown either way.
What protects against the remote threats is a key bound to the service,
which the standards allow as a single factor, but no source tests a key
alone, and the strongest evidence is for a key with user verification,
which is multi-factor.
Risk-based authentication and strong passwords are not substitutes.
Against a lost device a key alone protects nothing.
\Cref{h:necessary} is neither refuted nor established.

For the recommendation that prompted the question this means that
\enquote{use multi-factor authentication} is sound advice as a
correction of weak passwords, and that the advice which the evidence
supports better is narrower: use a factor that cannot be phished, and
protect it against loss of the device with something the device does
not hold.


\section{Future work}\label{sec:future}

Some of what this review left undone can be completed without new
research.
The 21 sources read only as abstracts should be read in full, first the
study of mandatory and optional use at banks
\parencite{Kamar2026TwoFactorBanks} and the experiments on push
approvals and overlays
\parencite{Jubur2021PushBypass,Mahdad2024FIDOLA}; four of the 21 are
freely available and were obtained too late to be read through.
Three documents could not be retrieved at all (\cref{sec:dropped}).
The four failed searches in IEEE Xplore should be repeated, and all
queries run on Semantic Scholar and DBLP.
A second, human screener should classify a sample of the records, to
measure what the model missed.

The questions that the evidence leaves open need new studies.
\begin{itemize}
  \item A measurement of compromise that separates the cases: accounts
    with a key alone, a key with user verification, and a password with
    a key, at one provider.
    If a password beside a key showed a measurable benefit, the answer
    to \cref{rq:attribution} would change; if a key alone did as well as
    the others, \cref{rq:necessary} would be answered with a no for the
    remote threats.
  \item A study of device loss in use: how often a stolen device leads
    to a compromised account, with and without user verification, and
    with what kind of PIN.
  \item A measurement with a comparison group of a strong, unique
    password against the same password with a second factor.
  \item Reports of incidents involving keys.
    The absence of any documented takeover of an account protected by a
    key is the weakest point of the conclusion on sufficiency: one such
    report would turn demonstrated possibility into observed fact.
\end{itemize}
This rapid review maps the ground.
The two questions on attribution and necessity are now precise enough
for the systematic review that \cref{sec:reviewtype} set aside.


\appendix

\section{Queries and records found}\label{sec:hitlists}

For each question this appendix gives the queries with the number of
records each database returned, and the list of the records that bear
on the question.
All searches were run on 12 November 2026 with the tool
\texttt{scholar}.
$M$ and $K$ are the term blocks of \cref{sec:queries}; $K'$ is $K$
without \q{single-factor} and with \kw{passkeys} added.
Web of Science received each query as \texttt{TS=(\dots)} and Scopus as
\texttt{TITLE-ABS-KEY(\dots)}, the other databases as written.
At most 40 records were taken per query and database, so 40 means
\enquote{at least 40}.
In the lists, a record marked \emph{cited} was read and is cited in the
report; the others carry the model's classification and its confidence.

% Author's note (verification and reproduction), all five searches:
%   scholar sessions : mfa-effect, mfa-attribution, mfa-sufficient,
%                      mfa-necessary, mfa-method
%   exports          : search/<session>.csv (classified session),
%                      search/<session>-full.tex (the tables input below)
%   searched         : 2026-11-12; providers openalex wos scopus ieee arxiv
%                      (s2: key rejected, HTTP 403 -- not searched; dblp:
%                      anti-bot page, HTTP 202 text/html -- not searched;
%                      ieee: daily quota of 200 requests exhausted after the
%                      planned queries -- four added queries failed)
%   classification   : scholar llm context + llm classify --no-examples
%                      -t supports-claim -t qualifies-claim
%                      -t adjacent-subtopic -t off-topic-false-hit;
%                      model anthropic/claude-sonnet-5-5 (mfa-effect:
%                      anthropic/claude-opus-4-8); human overrides 52
%                      (effect), 14 (attribution), 15 (sufficient), 9
%                      (necessary), written partly by editing the session
%                      files because `sessions decide` only appends tags
%   hit-list tables  : scholar sessions export <session> -f table
%                      --bearing-only --lang en --label hits-<name>
%                      --track "<title>" --theme TAG="name" ...
%                      (themes: effect measured-effect, bypass-attack;
%                      attribution local-malware, key-resists-remote,
%                      fallback-recovery, synced-passkeys,
%                      second-factor-weakness, formal-analysis; sufficient
%                      aitm, recovery, fidoatk, analysis; necessary RBA,
%                      FIDO, REC, STD, ATT; method review-types,
%                      shortcut-cost)
%   provenance       : mfa.bib, one block per use of a source

\subsection{Effect}

\begin{table}[h]
  \begin{sidecaption}{%
    Queries for \cref{rq:effect} and records returned by OpenAlex (OA;
    not open access), Web of Science (WoS), Scopus, IEEE Xplore (IEEE)
    and arXiv.
    S: worded to find support for the hypothesis, C: against it; an
    asterisk marks a query added after the first run.
  }[tab:q-effect]
    \footnotesize
    \begin{tabular}{@{}lp{0.42\linewidth}rrrrr@{}}
      \toprule
      & Keywords & OA & WoS & Scopus & IEEE & arXiv \\
      \midrule
      S1 & as in \cref{tab:planned} & 40 & 9 & 13 & 5 & 2 \\
      S2 & as in \cref{tab:planned} & 40 & 9 & 16 & 8 & 1 \\
      C1 & as in \cref{tab:planned} & 40 & 40 & 40 & 40 & 40 \\
      S3* & $M$ \kw{AND} (\q{compromised accounts} \kw{OR} \q{account
        compromise} \kw{OR} \q{account takeover} \kw{OR} \kw{hijacking}
        \kw{OR} \kw{hijacked}) \kw{AND} (\kw{reduces} \kw{OR}
        \kw{reduction} \kw{OR} \q{real-world} \kw{OR} \q{large-scale}
        \kw{OR} \kw{prevented} \kw{OR} \kw{blocks})
        & 40 & 25 & 29 & 18 & 3 \\
      S4* & $M$ \kw{AND} (\kw{adoption} \kw{OR} \kw{deployment} \kw{OR}
        \kw{enabled}) \kw{AND} (\kw{compromised} \kw{OR} \kw{breach}
        \kw{OR} \kw{hijacked} \kw{OR} \q{account takeover}) \kw{AND}
        (\kw{users} \kw{OR} \kw{accounts} \kw{OR} \kw{customers})
        & 40 & 39 & 40 & 40 & 10 \\
      \bottomrule
    \end{tabular}
  \end{sidecaption}
\end{table}
Outside the databases, seven web searches were made for measurements by
providers and for studies of self-selection among users of a second
factor; they led to the two provider statements cited and to the study
of Canadian banks, and found no study of self-selection.

% --- scholar LaTeX fragment --------------------------------------------
% Generated by: scholar sessions export -f table
% Session: mfa-effect
% \input this file from the body of a document whose
% preamble loads:
%   \usepackage{longtable}
%   \usepackage{booktabs}
%   \usepackage{array}
% ----------------------------------------------------------------------

\begingroup\footnotesize
\begin{longtable}{@{}%
  >{\raggedright\arraybackslash}p{0.44\textwidth}c%
  >{\raggedright\arraybackslash}p{0.13\textwidth}%
  >{\raggedright\arraybackslash}p{0.26\textwidth}@{}}
\caption{Records bearing on the claim, Effect of multi-factor authentication (9 cited, 0 supporting and 17 qualifying or contradicting it, out of 387 unique hits; the 290 adjacent and 69 false hits, 2 excluded although classified as bearing on the claim are not listed). The last column gives the reason for each inclusion; for machine-classified rows the model's confidence is given. Databases: OpenAlex (OA; not open access), IEEE Xplore (IEEE), Web of Science (WoS), Scopus, arXiv.}\label{tab:hits-effect}\\
\toprule
Title & Year & Provider & Reason \\
\midrule
\endfirsthead
\multicolumn{4}{@{}l}{\emph{\tablename~\thetable{} (cont.)}}\\
\toprule
Title & Year & Provider & Reason \\
\midrule
\endhead
\midrule \multicolumn{4}{r@{}}{\emph{continued on next page}}\\
\endfoot
\bottomrule
\endlastfoot
Defeating FIDO2/CTAP2/WebAuthn using browser in the middle and reflected cross site scripting & 2025 & WoS & \emph{cited}: Bypass or attack on MFA \\
Evaluating login challenges as a defense against account takeover & 2019 & Scopus & \emph{cited}: Measurement of effect \\
Evaluating the Effectiveness of Two-Factor Authentication (2FA) in Mitigating Account Takeover Fraud: A Natural Experimental Study on Canadian Banks & 2026 & Scopus & \emph{cited}: Measurement of effect \\
Hack for Hire: Exploring the emerging market for account hijacking & 2019 & Scopus & \emph{cited}: Bypass or attack on MFA \\
How effective is multifactor authentication at deterring cyberattacks? & 2023 & arXiv & \emph{cited}: Measurement of effect \\
Impersonation-as-a-Service: Characterizing the Emerging Criminal Infrastructure for User Impersonation at Scale & 2020 & arXiv & \emph{cited}: Bypass or attack on MFA \\
Internet fraud: The case of account takeover in online marketplace & 2019 & Scopus & \emph{cited}: Measurement of effect \\
Simple But Not Secure: An Empirical Security Analysis of Two-factor Authentication Systems & 2024 & arXiv & \emph{cited}: Bypass or attack on MFA \\
The Effectiveness of Two-Factor Authentication in Preventing Online Banking Fraud in Nigeria & 2026 & OA & \emph{cited}: Measurement of effect \\
Advanced Phishing Techniques: Analyzing Adversary-in-the-Middle and Browser-in-the-Browser Attacks in Modern Cybersecurity & 2025 & WoS & qualifies or contradicts the claim (0.90) \\
AI-Driven Identity Attacks: From Phishing to Session Hijacking & 2026 & WoS & qualifies or contradicts the claim (0.70) \\
All your face are belong to us & 2012 & OA & qualifies or contradicts the claim (0.85) \\
Can I still be hacked with 2FA enabled? & 2020 & OA & qualifies or contradicts the claim (0.65) \\
Defeating Multi-Factor Authentication with SIM Swapping & 2026 & OA & qualifies or contradicts the claim (0.90) \\
Faces in the Distorting Mirror & 2014 & OA & qualifies or contradicts the claim (0.85) \\
Impact Assessment of Password Reset PRMitM Attack with Two-Factor Authentication & 2019 & OA & qualifies or contradicts the claim (0.90) \\
Mobile Login Bridge: Subverting 2FA and Passwordless Authentication via Android Debug Bridge & 2024 & WoS & qualifies or contradicts the claim (0.85) \\
Phish in Sheep's Clothing: Exploring the Authentication Pitfalls of Browser Fingerprinting & 2022 & WoS, Scopus & qualifies or contradicts the claim (0.80) \\
PHISHING ATTACKS ON ENCRYPTED MESSENGERS: METHODS, RISKS AND PROTECTION RECOMMENDATIONS (USING THE EXAMPLE OF SIGNAL MESSENGER) & 2025 & OA & qualifies or contradicts the claim (0.70) \\
Security and Privacy Risks of Number Recycling at Mobile Carriers in the United States & 2021 & WoS & qualifies or contradicts the claim (0.85) \\
The 23andMe Data Breach: Analyzing Credential Stuffing Attacks, Security Vulnerabilities, and Mitigation Strategies & 2025 & arXiv & qualifies or contradicts the claim (0.80) \\
The 2FA Illusion: Uncovering Weak Links of Web Account Access in the Wild & 2025 & IEEE & qualifies or contradicts the claim (0.85) \\
The Adaptive Ecosystem of MaaS-Driven Cookie Theft: Dynamics, Anticipatory Analysis Concepts, and Proactive Defenses & 2025 & WoS & qualifies or contradicts the claim (0.80) \\
Thwarting Smartphone SMS Attacks at the Radio Interface Layer & 2023 & WoS & qualifies or contradicts the claim (0.75) \\
Two step verification system of highly secure social media: Possible to breach the security & 2017 & IEEE & qualifies or contradicts the claim (0.85) \\
Unveiling the Covert Vulnerabilities in Multi-Factor Authentication Protocols: A Systematic Review and Security Analysis & 2025 & WoS & qualifies or contradicts the claim (0.90) \\
\end{longtable}
\endgroup


\subsection{Attribution}

\begin{table}[h]
  \begin{sidecaption}{%
    Queries for \cref{rq:attribution} and records returned; columns and
    marks as in \cref{tab:q-effect}.
  }[tab:q-attribution]
    \footnotesize
    \begin{tabular}{@{}lp{0.42\linewidth}rrrrr@{}}
      \toprule
      & Keywords & OA & WoS & Scopus & IEEE & arXiv \\
      \midrule
      S1 & as in \cref{tab:planned} & 40 & 40 & 40 & 40 & 15 \\
      S2 & as in \cref{tab:planned} & 40 & 6 & 23 & 26 & 1 \\
      C1 & as in \cref{tab:planned} & 40 & 12 & 21 & 8 & 1 \\
      C2 & as in \cref{tab:planned} & 40 & 39 & 40 & 40 & 10 \\
      \bottomrule
    \end{tabular}
  \end{sidecaption}
\end{table}
Outside the databases, six web searches were made for standards and
guidance on how keys with and without user verification are classified;
they led to the guideline, the web standard, and the papers of the
United Kingdom's and the United States' agencies.
One standard, ISO/IEC 29115, is sold and was not obtained.

% --- scholar LaTeX fragment --------------------------------------------
% Generated by: scholar sessions export -f table
% Session: mfa-attribution
% \input this file from the body of a document whose
% preamble loads:
%   \usepackage{longtable}
%   \usepackage{booktabs}
%   \usepackage{array}
% ----------------------------------------------------------------------

\begingroup\footnotesize
\begin{longtable}{@{}%
  >{\raggedright\arraybackslash}p{0.44\textwidth}c%
  >{\raggedright\arraybackslash}p{0.13\textwidth}%
  >{\raggedright\arraybackslash}p{0.26\textwidth}@{}}
\caption{Records bearing on the claim, What each factor contributes (15 cited, 7 supporting and 10 qualifying or contradicting it, out of 289 unique hits; the 208 adjacent and 48 false hits, 1 excluded although classified as bearing on the claim are not listed). The last column gives the reason for each inclusion; for machine-classified rows the model's confidence is given. Databases: OpenAlex (OA; not open access), IEEE Xplore (IEEE), Web of Science (WoS), Scopus, arXiv.}\label{tab:hits-attribution}\\
\toprule
Title & Year & Provider & Reason \\
\midrule
\endfirsthead
\multicolumn{4}{@{}l}{\emph{\tablename~\thetable{} (cont.)}}\\
\toprule
Title & Year & Provider & Reason \\
\midrule
\endhead
\midrule \multicolumn{4}{r@{}}{\emph{continued on next page}}\\
\endfoot
\bottomrule
\endlastfoot
A Formal Analysis of the FIDO2 Protocols & 2022 & WoS & \emph{cited}: Formal analysis \\
A Security and Usability Analysis of Local Attacks Against FIDO2 & 2023 & arXiv, OA & \emph{cited}: Local attacks and malware \\
An Analysis of Attack Vectors Against FIDO2 Authentication & 2026 & arXiv & \emph{cited}: Local attacks and malware \\
An Extensive Formal Analysis of Multi-factor Authentication Protocols & 2021 & WoS & \emph{cited}: Local attacks and malware \\
Authentication Inconsistencies Across Online Services: A Multi-Scenario Security Analysis & 2025 & Scopus & \emph{cited}: Fallback and recovery \\
Beyond the Trust Boundary: A Critical Reassessment of the FIDO2 Threat Model & 2026 & OA & \emph{cited}: Local attacks and malware \\
Breaching Security Keys without Root: FIDO2 Deception Attacks via Overlays exploiting Limited Display Authenticators & 2024 & Scopus & \emph{cited}: Local attacks and malware \\
Bypassing Push-based Second Factor and Passwordless Authentication with Human-Indistinguishable Notifications & 2021 & WoS & \emph{cited}: Weak second factors \\
Device-Bound vs. Synced Credentials: A Comparative Evaluation of Passkey Authentication & 2025 & arXiv, OA & \emph{cited}: Synced passkeys \\
False Promises of Passwordless: Defeating Windows Hello through TPM Misuses & 2025 & Scopus & \emph{cited}: Local attacks and malware \\
Passwords and FIDO2 Are Meant To Be Secret: A Practical Secure Authentication Channel for Web Browsers & 2025 & OA & \emph{cited}: Local attacks and malware \\
Rogue key and impersonation attacks on FIDO2: From theory to practice & 2023 & OA & \emph{cited}: Local attacks and malware \\
Should We Rush to Implement Password-less Single Factor FIDO2 based Authentication? & 2020 & IEEE & \emph{cited}: A key resists remote attacks \\
SoK: A Comprehensive Evaluation of 2FA-based Schemes in the Face of Active Concurrent Attacks from User Terminal & 2023 & Scopus & \emph{cited}: Local attacks and malware \\
The 2FA Illusion: Uncovering Weak Links of Web Account Access in the Wild & 2025 & Scopus & \emph{cited}: Fallback and recovery \\
A Multi-Study Perspective on Phishing-Resistant Credentials and Employee Time-Based Resistance & 2026 & OA & supports the claim (0.45) \\
Comparing Multi-Factor Authentication Methods from SMS and TOTP to FIDO2/WebAuthn and Passkeys: A Qualitative Study of Practitioner Perspectives & 2026 & OA & supports the claim (0.35) \\
FIDO2 y WebAuthn para autenticación organizacional resistente al phishing: revisión sistemática & 2026 & OA & supports the claim (0.50) \\
Impact Assessment of Adversary-in-the-Middle Attack Evilginx Bypassing 2FA & 2025 & Scopus & supports the claim (0.50) \\
Passkeys and Fido2: Security, Usability and Deployment of Passwordless Authentication & 2026 & OA & supports the claim (0.55) \\
Simulation of Social Media Account Compromise Across User-Chosen Authentication Schemes & 2026 & Scopus & supports the claim (0.30) \\
Transition to Password-Free and Phishing-Resistant Authentication in Enterprise Information Systems: A FIDO2/Passkey-Based Standards-Based Qualitative Assessment and Phased Implementation Framework & 2026 & OA & supports the claim (0.50) \\
Beyond passwords: A review of the hidden risks in two-factor authentication & 2026 & Scopus & qualifies or contradicts the claim (0.50) \\
Defeating FIDO2/CTAP2/WebAuthn using browser in the middle and reflected cross site scripting & 2025 & OA & qualifies or contradicts the claim (0.60) \\
Evaluating the Security Posture of Real-World FIDO2 Deployments & 2023 & OA & qualifies or contradicts the claim (0.50) \\
Evaluation of Account Recovery Strategies with FIDO2-based Passwordless Authentication & 2021 & OA & qualifies or contradicts the claim (0.40) \\
HiPass: Hijacking CTAP in Passkey Authentication & 2025 & IEEE & qualifies or contradicts the claim (0.55) \\
Is FIDO2 Passwordless Authentication a Hype or for Real?: A Position Paper & 2022 & OA & qualifies or contradicts the claim (0.35) \\
Mobile Login Bridge: Subverting 2FA and Passwordless Authentication via Android Debug Bridge & 2024 & WoS & qualifies or contradicts the claim (0.55) \\
Passkey in the Wooden Horse: A Deep Dive into Trojan Based Compromise of Passwordless Credentials & 2026 & Scopus & qualifies or contradicts the claim (0.45) \\
Session Replication Attack Through QR Code Sniffing in Passkey CTAP Registration & 2024 & Scopus & qualifies or contradicts the claim (0.40) \\
USES OF WEBAUTHN PROTOCOL VULNERABILITIES TO OBTAIN UNSANCTIONED ACCESS & 2025 & OA & qualifies or contradicts the claim (0.55) \\
\end{longtable}
\endgroup


\subsection{Sufficiency}

\begin{table}[h]
  \begin{sidecaption}{%
    Queries for \cref{rq:sufficient} and records returned; columns and
    marks as in \cref{tab:q-effect}.
    Here S is worded to find attacks and C to find evidence of security.
    \enquote{--} means that the search failed: the daily quota of IEEE
    Xplore was used up.
  }[tab:q-sufficient]
    \footnotesize
    \begin{tabular}{@{}lp{0.42\linewidth}rrrrr@{}}
      \toprule
      & Keywords & OA & WoS & Scopus & IEEE & arXiv \\
      \midrule
      S1 & as in \cref{tab:planned} & 40 & 40 & 40 & 40 & 18 \\
      S2 & as in \cref{tab:planned} & 40 & 21 & 37 & 19 & 3 \\
      C1 & as in \cref{tab:planned} & 40 & 40 & 40 & 40 & 2 \\
      S3* & ($K'$) \kw{AND} (\kw{attack} \kw{OR} \kw{bypass} \kw{OR}
        \kw{downgrade} \kw{OR} \kw{vulnerability} \kw{OR} \kw{phishing}
        \kw{OR} \q{man-in-the-middle} \kw{OR}
        \q{adversary-in-the-middle} \kw{OR} \q{account recovery} \kw{OR}
        \q{session hijacking} \kw{OR} \kw{malware})
        & 40 & 40 & 40 & -- & 40 \\
      C2* & ($K'$) \kw{AND} (\q{formal verification} \kw{OR} \q{formal
        analysis} \kw{OR} \q{provable security} \kw{OR} \q{security
        proof} \kw{OR} \q{phishing-resistant} \kw{OR} \q{security
        analysis})
        & 40 & 40 & 40 & -- & 40 \\
      \bottomrule
    \end{tabular}
  \end{sidecaption}
\end{table}
Outside the databases, ten web searches were made for reviews of
incidents by public authorities and for reports by providers; they led
to the review of the Lapsus\$ intrusions, the guidance on
phishing-resistant authentication, and the three provider reports cited.
News articles were not used.

% --- scholar LaTeX fragment --------------------------------------------
% Generated by: scholar sessions export -f table
% Session: mfa-sufficient
% \input this file from the body of a document whose
% preamble loads:
%   \usepackage{longtable}
%   \usepackage{booktabs}
%   \usepackage{array}
% ----------------------------------------------------------------------

\begingroup\footnotesize
\begin{longtable}{@{}%
  >{\raggedright\arraybackslash}p{0.44\textwidth}c%
  >{\raggedright\arraybackslash}p{0.13\textwidth}%
  >{\raggedright\arraybackslash}p{0.26\textwidth}@{}}
\caption{Records bearing on the claim, Is multi-factor authentication sufficient (19 cited, 16 supporting and 7 qualifying or contradicting it, out of 508 unique hits; the 342 adjacent and 123 false hits, 1 excluded although classified as bearing on the claim are not listed). The last column gives the reason for each inclusion; for machine-classified rows the model's confidence is given. Databases: OpenAlex (OA; not open access), IEEE Xplore (IEEE), Web of Science (WoS), Scopus, arXiv.}\label{tab:hits-sufficient}\\
\toprule
Title & Year & Provider & Reason \\
\midrule
\endfirsthead
\multicolumn{4}{@{}l}{\emph{\tablename~\thetable{} (cont.)}}\\
\toprule
Title & Year & Provider & Reason \\
\midrule
\endhead
\midrule \multicolumn{4}{r@{}}{\emph{continued on next page}}\\
\endfoot
\bottomrule
\endlastfoot
A Formal Analysis of the FIDO2 Protocols & 2022 & WoS & \emph{cited}: Formal and analytical security analyses \\
A Security and Usability Analysis of Local Attacks Against FIDO2 & 2023 & arXiv & \emph{cited}: Attacks on FIDO2 and passkeys \\
Adventures in Recovery Land: Testing the Account Recovery of Popular Websites When the Second Factor is Lost & 2023 & WoS, Scopus & \emph{cited}: Recovery and fallback \\
An Analysis of Attack Vectors Against FIDO2 Authentication & 2026 & arXiv & \emph{cited}: Attacks on FIDO2 and passkeys \\
Beyond the Trust Boundary: A Critical Reassessment of the FIDO2 Threat Model & 2026 & arXiv & \emph{cited}: Attacks on FIDO2 and passkeys \\
CTRAPS: CTAP Client Impersonation and API Confusion on FIDO2 & 2025 & WoS & \emph{cited}: Attacks on FIDO2 and passkeys \\
Evaluating login challenges as a defense against account takeover & 2019 & Scopus & \emph{cited}: Recovery and fallback \\
Evaluation of Account Recovery Strategies with FIDO2-based Passwordless Authentication & 2021 & arXiv & \emph{cited}: Recovery and fallback \\
Excavating Vulnerabilities Lurking in Multi-Factor Authentication Protocols: A Systematic Security Analysis & 2024 & arXiv & \emph{cited}: Formal and analytical security analyses \\
False Promises of Passwordless: Defeating Windows Hello through TPM Misuses & 2025 & WoS & \emph{cited}: Attacks on FIDO2 and passkeys \\
FIDO2, CTAP 2.1, and WebAuthn 2: Provable Security and Post-Quantum Instantiation & 2023 & WoS & \emph{cited}: Formal and analytical security analyses \\
Impact Assessment of Adversary-in-the-Middle Attack Evilginx Bypassing 2FA & 2025 & Scopus & \emph{cited}: Phishing proxies and session theft \\
Impersonation-as-a-Service: Characterizing the Emerging Criminal Infrastructure for User Impersonation at Scale & 2020 & arXiv & \emph{cited}: Phishing proxies and session theft \\
Passkey in the Wooden Horse: A Deep Dive into Trojan Based Compromise of Passwordless Credentials & 2026 & Scopus & \emph{cited}: Attacks on FIDO2 and passkeys \\
Provable Security Analysis of FIDO2 & 2021 & WoS & \emph{cited}: Formal and analytical security analyses \\
Simple But Not Secure: An Empirical Security Analysis of Two-factor Authentication Systems & 2024 & arXiv & \emph{cited}: Phishing proxies and session theft \\
The Adaptive Ecosystem of MaaS-Driven Cookie Theft: Dynamics, Anticipatory Analysis Concepts, and Proactive Defenses & 2025 & Scopus & \emph{cited}: Phishing proxies and session theft \\
Traditional Multi Factor Authentication Under AiTM Phishing: A Controlled Campaign in Ecuador & 2027 & Scopus & \emph{cited}: Phishing proxies and session theft \\
Understanding Failures in Security Proofs of Multi-Factor Authentication for Mobile Devices & 2023 & WoS & \emph{cited}: Formal and analytical security analyses \\
``Abort the Login'': Understanding Phishing Susceptibility Through Warning Design in the Context of 2FA & 2026 & Scopus & supports the claim (0.75) \\
Account Recovery Methods for Two-Factor Authentication (2FA): An Exploratory Study & 2020 & OA & supports the claim (0.40) \\
AI-Driven Identity Attacks: From Phishing to Session Hijacking & 2026 & Scopus & supports the claim (0.55) \\
Analyzing 2FA Phishing Attacks and Their Prevention Techniques & 2022 & IEEE & supports the claim (0.70) \\
Authentication Inconsistencies Across Online Services: A Multi-Scenario Security Analysis & 2025 & Scopus & supports the claim (0.55) \\
Beyond Passwords: The Essence and Impact of Multi-Factor Authentication in Cybersecurity & 2025 & Scopus & supports the claim (0.60) \\
Defeating Multi-Factor Authentication with SIM Swapping & 2026 & OA & supports the claim (0.75) \\
Evaluating the Influence of Multi-Factor Authentication and Recovery Settings on the Security and Accessibility of User Accounts & 2024 & OA & supports the claim (0.70) \\
Exploiting DPAPI and Local State Decryption for Web Cookie Session Theft in Cross-Device Chrome Migrations & 2025 & Scopus & supports the claim (0.85) \\
Generative AI Deepfakes: Real-World Threats to Facial Biometric Authentication & 2025 & WoS & supports the claim (0.80) \\
Multi-Factor Authentication Bypass - Attack Techniques and :Protection Mechanisms & 2026 & OA & supports the claim (0.90) \\
Pixnapping: Bringing Pixel Stealing out of the Stone Age & 2025 & WoS & supports the claim (0.45) \\
SoK: Web Authentication and Recovery in the Age of End-to-End Encryption & 2025 & OA & supports the claim (0.50) \\
The Mobile Paradox: How Interface Friction Can Improve 2FA Security & 2026 & Scopus & supports the claim (0.55) \\
User-centric security analysis of MitID: The Danish passwordless digital identity solution & 2023 & WoS & supports the claim (0.55) \\
Your Code is 0000: An Analysis of the Disposable Phone Numbers Ecosystem & 2023 & arXiv & supports the claim (0.80) \\
Assessment of multi factor authentication techniques for mitigating phishing attack & 2026 & OA & qualifies or contradicts the claim (0.70) \\
Comparative Security Analysis of FIDO2/WebAuthn versus Traditional Multi-Factor Authentication: A Systematic Review & 2026 & OA & qualifies or contradicts the claim (0.60) \\
Device-Bound vs. Synced Credentials: A Comparative Evaluation of Passkey Authentication & 2025 & OA & qualifies or contradicts the claim (0.70) \\
FIDO2 y WebAuthn para autenticación organizacional resistente al phishing: revisión sistemática & 2026 & OA & qualifies or contradicts the claim (0.55) \\
From Passwords to Passkeys. Design, implementation and security evaluation of WebAuthn-Based passwordless authentication & 2026 & OA & qualifies or contradicts the claim (0.45) \\
Passkeys and Fido2: Security, Usability and Deployment of Passwordless Authentication & 2026 & OA & qualifies or contradicts the claim (0.70) \\
Should We Rush to Implement Password-less Single Factor FIDO2 based Authentication? & 2020 & OA & qualifies or contradicts the claim (0.50) \\
\end{longtable}
\endgroup


\subsection{Necessity}

\begin{table}[h]
  \begin{sidecaption}{%
    Queries for \cref{rq:necessary} and records returned; columns and
    marks as in \cref{tab:q-effect,tab:q-sufficient}.
    Here S is worded to find alternatives to several factors and C to
    find evidence that several are needed.
  }[tab:q-necessary]
    \footnotesize
    \begin{tabular}{@{}lp{0.42\linewidth}rrrrr@{}}
      \toprule
      & Keywords & OA & WoS & Scopus & IEEE & arXiv \\
      \midrule
      S1 & as in \cref{tab:planned} & 40 & 30 & 28 & 27 & 2 \\
      S2 & as in \cref{tab:planned} & 40 & 40 & 40 & 40 & 12 \\
      C1 & as in \cref{tab:planned} & 40 & 40 & 40 & 40 & 40 \\
      C2 & as in \cref{tab:planned} & 40 & 27 & 36 & 14 & 1 \\
      S3* & ($K$) \kw{AND} (\kw{phishing} \kw{OR} \q{account takeover}
        \kw{OR} \q{credential stuffing}) \kw{AND} (\kw{evaluation}
        \kw{OR} \kw{measurement} \kw{OR} \q{security analysis} \kw{OR}
        \q{formal analysis})
        & 40 & 12 & 20 & -- & 8 \\
      C3* & $M$ \kw{AND} (\q{assurance level} \kw{OR}
        \q{phishing-resistant} \kw{OR} \q{NIST SP 800-63B} \kw{OR}
        \q{strong customer authentication})
        & 40 & 13 & 21 & -- & 1 \\
      \bottomrule
    \end{tabular}
  \end{sidecaption}
\end{table}
One search of Web of Science timed out and was repeated unchanged.
Outside the databases, six web searches were made for standards,
regulations and guidance, and two for measured comparisons of strong
passwords and of keys with multi-factor authentication; the latter two
returned only vendors' pages, which are not cited.

% --- scholar LaTeX fragment --------------------------------------------
% Generated by: scholar sessions export -f table
% Session: mfa-necessary
% \input this file from the body of a document whose
% preamble loads:
%   \usepackage{longtable}
%   \usepackage{booktabs}
%   \usepackage{array}
% ----------------------------------------------------------------------

\begingroup\footnotesize
\begin{longtable}{@{}%
  >{\raggedright\arraybackslash}p{0.44\textwidth}c%
  >{\raggedright\arraybackslash}p{0.13\textwidth}%
  >{\raggedright\arraybackslash}p{0.26\textwidth}@{}}
\caption{Records bearing on the claim, Is multi-factor authentication necessary (21 cited, 11 supporting and 27 qualifying or contradicting it, out of 2128 unique hits; the 1641 adjacent and 428 false hits are not listed). The last column gives the reason for each inclusion; for machine-classified rows the model's confidence is given. Databases: OpenAlex (OA; not open access), IEEE Xplore (IEEE), Web of Science (WoS), Scopus, arXiv.}\label{tab:hits-necessary}\\
\toprule
Title & Year & Provider & Reason \\
\midrule
\endfirsthead
\multicolumn{4}{@{}l}{\emph{\tablename~\thetable{} (cont.)}}\\
\toprule
Title & Year & Provider & Reason \\
\midrule
\endhead
\midrule \multicolumn{4}{r@{}}{\emph{continued on next page}}\\
\endfoot
\bottomrule
\endlastfoot
A Study of Multi-Factor and Risk-Based Authentication Availability & 2023 & WoS, Scopus & \emph{cited}: Risk-based authentication \\
An Analysis of Attack Vectors Against FIDO2 Authentication & 2026 & OA & \emph{cited}: Cryptographic keys and passkeys \\
An Extensive Formal Analysis of Multi-factor Authentication Protocols & 2021 & Scopus, OA & \emph{cited}: Attacks on and measurements of MFA \\
Best Practices for Privileged User PIV Authentication & 2016 & OA & \emph{cited}: Standards and requirements \\
Breaching Security Keys without Root: FIDO2 Deception Attacks via Overlays exploiting Limited Display Authenticators & 2024 & OA & \emph{cited}: Cryptographic keys and passkeys \\
Defeating FIDO2/CTAP2/WebAuthn using browser in the middle and reflected cross site scripting & 2025 & OA & \emph{cited}: Cryptographic keys and passkeys \\
Device-Bound vs. Synced Credentials: A Comparative Evaluation of Passkey Authentication & 2025 & OA & \emph{cited}: Cryptographic keys and passkeys \\
Evaluating the Effectiveness of Two-Factor Authentication (2FA) in Mitigating Account Takeover Fraud: A Natural Experimental Study on Canadian Banks & 2026 & OA & \emph{cited}: Attacks on and measurements of MFA \\
Evaluation of account recovery strategies with FIDO2-based passwordless authentication & 2021 & Scopus, OA & \emph{cited}: Account recovery and fallback \\
Evaluation of Real-World Risk-Based Authentication at Online Services Revisited: Complexity Wins & 2023 & OA & \emph{cited}: Risk-based authentication \\
Evidence-based cybersecurity policy? A meta-review of security control effectiveness & 2023 & OA & \emph{cited}: Standards and requirements \\
False Promises of Passwordless: Defeating Windows Hello through TPM Misuses & 2025 & OA & \emph{cited}: Cryptographic keys and passkeys \\
Impersonation-as-a-Service: Characterizing the Emerging Criminal Infrastructure for User Impersonation at Scale & 2020 & OA & \emph{cited}: Attacks on and measurements of MFA \\
Is It Really You Who Forgot the Password? When Account Recovery Meets Risk-Based Authentication & 2024 & OA & \emph{cited}: Risk-based authentication \\
Is this really you? An empirical study on risk-based authentication applied in the wild & 2019 & OA & \emph{cited}: Risk-based authentication \\
Mobile device fingerprinting considered harmful for risk-based authentication & 2015 & OA & \emph{cited}: Risk-based authentication \\
Multi-Factor Authentication for Criminal Justice Information Systems: & 2024 & OA & \emph{cited}: Standards and requirements \\
NIST SP 800-63B-4: & 2025 & OA & \emph{cited}: Standards and requirements \\
Phish in Sheep's Clothing: Exploring the Authentication Pitfalls of Browser Fingerprinting & 2022 & WoS, Scopus & \emph{cited}: Risk-based authentication \\
Pump Up Password Security! Evaluating and Enhancing Risk-Based Authentication on a Real-World Large-Scale Online Service & 2022 & OA & \emph{cited}: Risk-based authentication \\
The 23andMe Data Breach: Analyzing Credential Stuffing Attacks, Security Vulnerabilities, and Mitigation Strategies & 2025 & OA & \emph{cited}: Attacks on and measurements of MFA \\
Comparative Analysis of User Authentication Technologies for On-Premises Cloud Systems with Key Derivation: Face Biometrics, PIN and Fuzzy Extractors & 2026 & OA & supports the claim (0.50) \\
Comparative Security Analysis of FIDO2/WebAuthn versus Traditional Multi-Factor Authentication: A Systematic Review & 2026 & OA & supports the claim (0.60) \\
Comparing Multi-Factor Authentication Methods from SMS and TOTP to FIDO2/WebAuthn and Passkeys: A Qualitative Study of Practitioner Perspectives & 2026 & OA & supports the claim (0.40) \\
Enhancing Digital Wallet Security: A Systematic Comparison of Passwordless and Risk-Based Authentication Approaches & 2026 & OA & supports the claim (0.35) \\
Enhancing Identity and Access Management in the U.S. Navy Via Migration to More Modern Standards of Authentication & 2020 & OA & supports the claim (0.35) \\
FIDO2 y WebAuthn para autenticación organizacional resistente al phishing: revisión sistemática & 2026 & OA & supports the claim (0.50) \\
Passkeys and Fido2: Security, Usability and Deployment of Passwordless Authentication & 2026 & OA & supports the claim (0.60) \\
Read This before You Take Multi-Factor Plunge: Regulators Say Multi-Factor Authentication Is an Option Not a Requirement & 2006 & OA & supports the claim (0.50) \\
Resilience of the FIDO2/WebAuthn Protocol and Hardware Security Keys (YubiKey) Against Adversary-in-the-Middle (AitM) Phishing Attacks Compared to RFC 6238 TOTP: An Empirical Study and Architectural Framework & 2026 & OA & supports the claim (0.80) \\
Transition to Password-Free and Phishing-Resistant Authentication in Enterprise Information Systems: A FIDO2/Passkey-Based Standards-Based Qualitative Assessment and Phased Implementation Framework & 2026 & OA & supports the claim (0.50) \\
What's in Score for Website Users: A Data-Driven Long-Term Study on Risk-Based Authentication Characteristics & 2021 & OA & supports the claim (0.70) \\
A Comparative Analysis of NIST SP 800-63B-4 2pd(2024) and Domestic Password Policies & 2024 & OA & qualifies or contradicts the claim (0.45) \\
A Longitudinal Study on Web-Sites Password Management (in)Security: Evidence and Remedies & 2020 & OA & qualifies or contradicts the claim (0.45) \\
A systematic review of multi-factor authentication in digital payment systems: NIST standards alignment and industry implementation analysis & 2025 & OA & qualifies or contradicts the claim (0.50) \\
AI-Driven Identity Attacks: From Phishing to Session Hijacking & 2026 & OA & qualifies or contradicts the claim (0.50) \\
An In-Depth Analysis of Password Managers and Two-Factor Authentication Tools & 2025 & WoS & qualifies or contradicts the claim (0.45) \\
Assessment of multi factor authentication techniques for mitigating phishing attack & 2026 & OA & qualifies or contradicts the claim (0.85) \\
Authentication Inconsistencies Across Online Services: A Multi-Scenario Security Analysis & 2025 & OA & qualifies or contradicts the claim (0.50) \\
Beyond Digital Knowledge: The Dominant Roles of User Vigilance and Two-Factor Authentication in E-Commerce Phishing Susceptibility & 2026 & OA & qualifies or contradicts the claim (0.40) \\
Beyond Passwords: The Essence and Impact of Multi-Factor Authentication in Cybersecurity & 2025 & OA & qualifies or contradicts the claim (0.50) \\
Data security and multi-factor authentication: Analysis of requirements under EU law and in selected EU Member States & 2016 & OA & qualifies or contradicts the claim (0.75) \\
E-Banking Security Study---10 Years Later & 2022 & OA & qualifies or contradicts the claim (0.80) \\
Eadem, Sed Aliter: Investigating the Instability of Browser Fingerprinting & 2026 & OA & qualifies or contradicts the claim (0.60) \\
Evaluation of System Access Security in The Implementation of Multi-Factor Authentication (MFA) in Educational Institutions & 2024 & OA & qualifies or contradicts the claim (0.50) \\
EVALUATION OF TWO-FACTOR AUTHENTICATION METHOD IN IMPROVING AUTHENTICATION SECURITY ON SSL VPN (Case Study of PT. Kanmo Group) & 2025 & OA & qualifies or contradicts the claim (0.60) \\
Formal modelling and verifying eIDAS multi-factor authentication with interface-based threat analysis & 2026 & OA & qualifies or contradicts the claim (0.55) \\
Formal Verification of Peer-Assisted FIDO2 Passkey Recovery Protocol with Tamarin & 2026 & Scopus & qualifies or contradicts the claim (0.40) \\
From Passwords to Passkeys. Design, implementation and security evaluation of WebAuthn-Based passwordless authentication & 2026 & OA & qualifies or contradicts the claim (0.50) \\
Ich sehe was, das du nicht siehst - Die Realität von Mobilebanking zwischen allgemeinen und rechtlichen Anforderungen & 2018 & OA & qualifies or contradicts the claim (0.55) \\
Impact Assessment of Adversary-in-the-Middle Attack Evilginx Bypassing 2FA & 2025 & OA & qualifies or contradicts the claim (0.75) \\
Is FIDO2 Passwordless Authentication a Hype or for Real?: A Position Paper & 2022 & OA & qualifies or contradicts the claim (0.50) \\
Multifactor authentication for e-commerce & 2019 & OA & qualifies or contradicts the claim (0.40) \\
Practical Acoustic Eavesdropping On Typed Passphrases & 2025 & OA & qualifies or contradicts the claim (0.70) \\
Revisiting Authentication Security: Threat Models, Limitations of Single-Factor Systems, and the Case for Multi-Factor Authentication & 2026 & OA & qualifies or contradicts the claim (0.50) \\
The Second-Factor Authentication System at CERN & 2024 & OA & qualifies or contradicts the claim (0.80) \\
Traditional Multi Factor Authentication Under AiTM Phishing: A Controlled Campaign in Ecuador & 2027 & OA & qualifies or contradicts the claim (0.75) \\
Two-Factor Authentication: A Solution to Times Past or Present? The Debate Surrounding the Gramm-Leach-Bliley Security Safeguards Rule and the Methods of Risk Assessment and Compliance & 2006 & OA & qualifies or contradicts the claim (0.50) \\
Unveiling the Covert Vulnerabilities in Multi-Factor Authentication Protocols: A Systematic Review and Security Analysis & 2025 & OA & qualifies or contradicts the claim (0.50) \\
\end{longtable}
\endgroup


\subsection{Types of review}\label{sec:hits-method}

The claims of \cref{sec:reviewtype} about types of review and about what
the shortcuts of a rapid review cost were backed by a search of their
own, with one query in each direction.
The course reading cited there was retrieved as a known document and
is not among the records.
\begin{table}[h]
  \begin{sidecaption}{%
    Queries on types of review and records returned; columns as in
    \cref{tab:q-effect}.
    S: worded to find definitions of the types, C: to find what
    shortcuts cost or do not cost.
  }[tab:q-method]
    \footnotesize
    \begin{tabular}{@{}lp{0.42\linewidth}rrrrr@{}}
      \toprule
      & Keywords & OA & WoS & Scopus & IEEE & arXiv \\
      \midrule
      S1 & \q{rapid review} \kw{AND} \q{scoping review} \kw{AND}
        \q{systematic review} \kw{AND} (\kw{typology} \kw{OR} \q{review
        types} \kw{OR} \kw{definitions} \kw{OR} \kw{methodologies})
        & 40 & 19 & 40 & 1 & 0 \\
      C1 & \q{rapid review} \kw{AND} (\q{single reviewer} \kw{OR}
        \q{single screening} \kw{OR} \kw{shortcuts} \kw{OR}
        \q{streamlined}) \kw{AND} (\kw{bias} \kw{OR} \kw{accuracy}
        \kw{OR} \q{missed studies} \kw{OR} \kw{limitations})
        & 40 & 35 & 40 & 1 & 0 \\
      \bottomrule
    \end{tabular}
  \end{sidecaption}
\end{table}

% --- scholar LaTeX fragment --------------------------------------------
% Generated by: scholar sessions export -f table
% Session: mfa-method
% \input this file from the body of a document whose
% preamble loads:
%   \usepackage{longtable}
%   \usepackage{booktabs}
%   \usepackage{array}
% ----------------------------------------------------------------------

\begingroup\footnotesize
\begin{longtable}{@{}%
  >{\raggedright\arraybackslash}p{0.44\textwidth}c%
  >{\raggedright\arraybackslash}p{0.13\textwidth}%
  >{\raggedright\arraybackslash}p{0.26\textwidth}@{}}
\caption{Records bearing on the claim, Types of review and the cost of shortcuts (6 cited, 54 supporting and 3 qualifying or contradicting it, out of 133 unique hits; the 30 adjacent and 40 false hits are not listed). The last column gives the reason for each inclusion; for machine-classified rows the model's confidence is given. Databases: OpenAlex (OA; not open access), IEEE Xplore (IEEE), Web of Science (WoS), Scopus, arXiv.}\label{tab:hits-method}\\
\toprule
Title & Year & Provider & Reason \\
\midrule
\endfirsthead
\multicolumn{4}{@{}l}{\emph{\tablename~\thetable{} (cont.)}}\\
\toprule
Title & Year & Provider & Reason \\
\midrule
\endhead
\midrule \multicolumn{4}{r@{}}{\emph{continued on next page}}\\
\endfoot
\bottomrule
\endlastfoot
A scoping review of rapid review methods & 2015 & WoS & \emph{cited}: Types of review \\
Excluding non-English publications from evidence-syntheses did not change conclusions: a meta-epidemiological study & 2020 & Scopus & \emph{cited}: What shortcuts cost \\
Rapid Reviews in Software Engineering & 2020 & OA & \emph{cited}: Types of review \\
Rapid reviews may produce different results to systematic reviews: a meta-epidemiological study & 2018 & OA & \emph{cited}: What shortcuts cost \\
Single screening versus conventional double screening for study selection in systematic reviews: A methodological systematic review & 2019 & Scopus & \emph{cited}: What shortcuts cost \\
Systematic review or scoping review? Guidance for authors when choosing between a systematic or scoping review approach & 2018 & OA & \emph{cited}: Types of review \\
Advancing knowledge of rapid reviews: an analysis of results, conclusions and recommendations from published review articles examining rapid reviews & 2015 & OA & supports the claim (0.75) \\
Advancing the Use of Knowledge Synthesis to Inform Policy and Decision Making in Agri-food Public Health & 2014 & OA & supports the claim (0.45) \\
All in the Family: systematic reviews, rapid reviews, scoping reviews, realist reviews, and more & 2015 & OA & supports the claim (0.70) \\
An Analysis of Literature Reviews in the Context of Healthcare Program Assessment & 2021 & OA & supports the claim (0.75) \\
An Overview of Review Types: Foundational Support & 2024 & Scopus & supports the claim (0.60) \\
Characteristics and recovery methods of studies falsely excluded during literature screening---a systematic review & 2022 & Scopus & supports the claim (0.45) \\
Characteristics of studies falsely excluded during single-reviewer abstract screening: a meta-epidemiological analysis & 2026 & Scopus & supports the claim (0.70) \\
Choose your shortcuts wisely: COVID-19 rapid reviews of traditional, complementary and integrative medicine & 2020 & Scopus & supports the claim (0.45) \\
Cochrane Rapid Reviews Methods Group offers evidence-informed guidance to conduct rapid reviews & 2020 & OA & supports the claim (0.50) \\
Combining abbreviated literature searches with single-reviewer screening: Three case studies of rapid reviews & 2020 & Scopus & supports the claim (0.60) \\
Come sintetizzare i risultati degli studi biomedici: una rassegna metodologica & 2026 & OA & supports the claim (0.70) \\
Defining rapid reviews: a systematic scoping review and thematic analysis of definitions and defining characteristics of rapid reviews & 2021 & Scopus & supports the claim (0.80) \\
Designing a rapid response program to support evidence-informed decision-making in the Americas region: using the best available evidence and case studies & 2016 & WoS & supports the claim (0.45) \\
Evaluating the accuracy and impact of methodological trade-offs used in rapid reviews on results and certainty of the evidence: meta-research studies in systematic reviews of chronic low back pain & 2026 & OA & supports the claim (0.55) \\
Evaluating the accuracy of full text and serial screening by one vs two independent reviewers: part 2 of a meta-research study of rapid review methodological trade-offs & 2026 & Scopus & supports the claim (0.60) \\
Evaluating the accuracy of screening titles and abstracts for progressing to full-text screening by one vs two independent reviewers: part 1 of a meta-research study of rapid review methodological trade-offs & 2026 & Scopus & supports the claim (0.60) \\
Evaluating the impact of methodological trade-offs on meta-analyses and certainty of the evidence according to GRADE: part 5 of a meta-research study of rapid reviews & 2026 & Scopus & supports the claim (0.55) \\
Evidence summaries: the evolution of a rapid review approach & 2012 & OA & supports the claim (0.50) \\
Evidenzsynthesen in Public Health: ein Überblick & 2022 & OA & supports the claim (0.65) \\
Expediting systematic reviews: methods and implications of rapid reviews & 2010 & WoS & supports the claim (0.85) \\
Few evaluative studies exist examining rapid review methodology across stages of conduct: a systematic scoping review & 2020 & Scopus & supports the claim (0.65) \\
Implications of applying methodological shortcuts to expedite systematic reviews: three case studies using systematic reviews from agri-food public health & 2016 & WoS & supports the claim (0.90) \\
Key concepts in rapid reviews: an overview & 2024 & OA & supports the claim (0.70) \\
Knowledge Syntheses in Medical Education: Demystifying Scoping Reviews & 2016 & OA & supports the claim (0.70) \\
Literature Review and Synthesis Typology & 2022 & OA & supports the claim (0.60) \\
Meeting the review family: exploring review types and associated information retrieval requirements & 2019 & OA & supports the claim (0.70) \\
One question, many reviews: practical typology and decision guide for psychiatrists and psychologists & 2026 & Scopus & supports the claim (0.65) \\
Over 1000 terms have been used to describe evidence synthesis: a scoping review & 2025 & OA & supports the claim (0.50) \\
Overview of evidence synthesis types and modes & 2025 & Scopus & supports the claim (0.65) \\
Paving the way for mature secondary research: the seven types of literature review & 2022 & Scopus & supports the claim (0.60) \\
PP109 Which Review Is Right For You? Choosing A Review Methodology & 2024 & OA & supports the claim (0.80) \\
Quality of conduct and reporting in rapid reviews: an exploration of compliance with PRISMA and AMSTAR guidelines & 2016 & OA & supports the claim (0.60) \\
Rapid review: A review of methods and recommendations based on current evidence & 2024 & Scopus & supports the claim (0.60) \\
Rapid reviews methods series: assessing the appropriateness of conducting a rapid review & 2024 & OA & supports the claim (0.70) \\
Rapid reviews: the pros and cons of an accelerated review process & 2021 & OA & supports the claim (0.70) \\
Redefining rapid reviews: a flexible framework for restricted systematic reviews & 2018 & OA & supports the claim (0.65) \\
Reviewing the animal literature: how to describe and choose between different types of literature reviews & 2021 & WoS & supports the claim (0.70) \\
Risk of Bias in Rapid Reviews: Omission, Adaptation and Transparency & 2026 & OA & supports the claim (0.60) \\
Same family, different species: methodological conduct and quality varies according to purpose for five types of knowledge synthesis & 2017 & OA & supports the claim (0.80) \\
Scoping studies: towards a methodological framework & 2005 & OA & supports the claim (0.80) \\
Single-reviewer abstract screening missed 13 percent of relevant studies: a crowd-based, randomized controlled trial & 2020 & Scopus & supports the claim (0.90) \\
Syntheses Synthesized: A Look Back at Grant and Booth's Review Typology & 2022 & Scopus & supports the claim (0.65) \\
Systematic and Other Reviews: Criteria and Complexities & 2021 & Scopus & supports the claim (0.50) \\
Systematic Review vs Scoping Review: Scoping Review with SAIMSARA & 2026 & OA & supports the claim (0.75) \\
Systematic Versus Rapid Versus Scoping Reviews & 2022 & Scopus & supports the claim (0.80) \\
Systematic, narrative, scoping or rapid reviews - when, why and how to choose the best option in pharmacy education & 2026 & Scopus & supports the claim (0.70) \\
Types of literature reviews in health: from narrative review to systematic review with meta-analysis & 2026 & OA & supports the claim (0.70) \\
Typology and Methodology of Reviews in Health Research & 2022 & Scopus & supports the claim (0.65) \\
Unbiased Approaches for Reviewing Entomology Literature: A Systematized Review & 2021 & WoS & supports the claim (0.45) \\
Updated recommendations for the Cochrane rapid review methods guidance for rapid reviews of effectiveness & 2024 & OA & supports the claim (0.55) \\
Use of evidence-based practice in an aid organisation & 2014 & OA & supports the claim (0.45) \\
What are the best methodologies for rapid reviews of the research evidence for evidence-informed decision making in health policy and practice: a rapid review & 2016 & WoS & supports the claim (0.65) \\
What are the best methods for rapid reviews of the research evidence? A systematic review of reviews and primary studies & 2024 & Scopus & supports the claim (0.65) \\
When Is The Use Of A Systematic Literature Review Appropriate? A Comparison Of Systematic, Rapid, And Scoping Reviews And Their Application To The Hta Process & 2016 & OA & supports the claim (0.80) \\
A retrospective comparison of systematic reviews with same-topic rapid reviews & 2017 & OA & qualifies or contradicts the claim (0.50) \\
Comparison of a full systematic review versus rapid review approaches to assess a newborn screening test for tyrosinemia type 1 & 2017 & Scopus & qualifies or contradicts the claim (0.45) \\
Evaluating the proportion of retrieved citations and relevant studies when database selection is restricted to MEDLINE, CENTRAL, and Embase: part 4 of a meta-research study of rapid review methodological trade-offs & 2026 & Scopus & qualifies or contradicts the claim (0.45) \\
\end{longtable}
\endgroup



\section{Sources dropped and documents not obtained}\label{sec:dropped}

Three sources that a search run had cited were dropped in the second
check.
\begin{itemize}
  \item A report of laboratory trials of keys against relayed logins,
    published by its single author in a repository without review, of
    which only the abstract could be read (DOI 10.5281/zenodo.23955576).
    The claims it supported are carried by other sources.
  \item A trade-press article of 2006 about guidance from United States
    banking regulators, known only from a garbled abstract.
    It reports the guidance at second hand; the guidance itself could
    not be retrieved.
  \item Two studies of multi-factor authentication at a university and
    in banking in Ghana (DOIs 10.56201/ijcsmt.vol.12.no1.2026.pg144.154
    and 10.37602/ijssmr.2025.8106), whose abstracts state no size of
    effect.
\end{itemize}
Three documents that the questions call for could not be retrieved,
because the issuing site refused automated retrieval or the text is
sold.
\begin{itemize}
  \item The European Union's regulatory technical standard on strong
    customer authentication, Delegated Regulation (EU) 2018/389.
  \item The guidance of 2005 on authentication in internet banking by
    the United States Federal Financial Institutions Examination
    Council.
  \item ISO/IEC 29115, the international standard on assurance of
    entity authentication.
\end{itemize}
